\chapter{Probability Guarantees and Positive Realization}\label{sec:probability}

\section{From a linear representation to a probability law}
Finite Hankel rank describes the dimension of a linear realization
\citep{carlyle1971,fliess1974}. It
does not impose nonnegative transitions, almost-sure termination, or
robust predictions after conditioning on a rare prefix. These constraints
are distinct and must enter a probabilistic reduction theorem separately.
Section~\ref{sec:hankel} treats the operator approximation problem.
Here we ask when spectral information or an acquired dictionary can
certify a proper approximation of a finite-word law.

Let $\Sigma$ have $d\ge2$ letters and let $p$ be a probability law on
completed words, including the empty word. All logarithms in this section
are natural. The word-Hankel operator is $H_p[u,v]=p(uv)$ and
$t_p(T)=\Pr_p\{|W|>T\}$ for integer $T\ge0$. We assume
$p(\varepsilon)\ge\mu$ with $0<\mu<1$ when using the empty pivot in the spectral
bridge. Let $\mathcal Q_N^+$ be the proper nonnegative automata with
at most $N$ states, and write
\[
 \mathcal K(p)=\inf_{q\in\mathcal Q_1^+}D_{\rm KL}(p\|q).
\]
A one-state law consists of an independent stop/letter event distribution.
Its information projection is explicit: with $M=\mathbb E_p|W|$,
letter-frequency law $\pi$, and
$g(M)=(M+1)\log(M+1)-M\log M$,
\[
 \mathcal K(p)=g(M)+M H(\pi)-H(p)
\]
when the terms are finite, with the conventions in
Proposition~\ref{II:U:prop:iprojection}. This classical calculation fixes
the comparison class; the new issue is the sharp dependence on a small
second Hankel singular value.

\section{Sharp spectral-to-probability regimes}
\begin{theorem}[Three length regimes]
Fix $0<\mu<1$, $C>0$, $M_s>0$, $M_0>0$, and the alphabet size; use the indicated parameters in each class.
\begin{enumerate}
\item If $t_p(T)\le C\rho^{T+1}$ for all integers $T\ge0$, where
$0<\rho<1$, then
\[
 \sup_{\substack{p(\varepsilon)\ge\mu,\ t_p(T)\le C\rho^{T+1}\\
                   \sigma_2(H_p)\le\epsilon}}\mathcal K(p)
 =\Theta\bigl(\epsilon^\gamma\log^2(1/\epsilon)\bigr),\qquad
 \gamma=\frac{\log(1/\rho)}{\log(1/\rho)+\tfrac12\log d}.
\]
\item If $s>1$ and $\mathbb E|W|^s\le M_s$, the corresponding supremum is
\[
 \Theta\bigl((\log(1/\epsilon))^{1-s}\log\log(1/\epsilon)\bigr).
\]
\item Under only $\mathbb E|W|\le M_0$, that supremum equals
$g(M_0)+M_0\log d$ for every $\epsilon>0$, but equals zero at
$\epsilon=0$.
\end{enumerate}
The asymptotic bounds use sufficiently small
$0<\epsilon\le\mu/2$ and integer truncation lengths. Constants depend
on the fixed class parameters; they are not uniform as $s\downarrow1$.
\end{theorem}
The exact statements and proofs are
Theorems~\ref{II:U:thm:kl-exp}, \ref{II:U:thm:kl-moment} and
\ref{II:U:thm:kl-mean}.

The proof first converts the rank-one spectral defect into a proper law
$q$ through the nonvanishing empty pivot. For every integer $T\ge0$,
\[
 \operatorname{TV}(p,q)
 \le\frac{dC_\mu}{\mu}\epsilon\sqrt{\sum_{j=0}^Td^j}
       +\frac{t_p(T)}{\mu^2},\qquad
 C_\mu=(1+1/\mu)(1+2/\mu).
\]
The first term is a finite-head Cauchy--Schwarz estimate; the second uses
positivity and the source tail. Length-tail information then controls
entropy and the cross-entropy loss after smoothing the event law.
Explicit integer choices of $T$ and $0<\Delta\le1/2$ are given in the
proofs. Smoothing is a step in the bound, not a claim that an arbitrary
spectral truncation is positive.

Sharp lower bounds use laws concentrating a small probability on a long
length shell. For
$p_T=(1-\delta)\delta_\varepsilon+\delta\operatorname{Unif}(\Sigma^T)$,
$T\ge2$, the shell-basis singular decomposition gives
$\sigma_2(H_{p_T})=\delta d^{-T/2}$ exactly, while
$\mathcal K(p_T)=g(\delta T)-h_2(\delta)$.
The choices of $\delta$ corresponding to exponential and moment tails
produce the stated rates. Under a mean bound, a single long word with
balanced letter counts supplies the discontinuity. A separate argument
in Theorem~\ref{II:U:thm:kl-fixed-states} shows that allowing any fixed
number of positive states does not restore a universal spectral-to-KL
continuity theorem under the hypotheses stated there.

\section{A constructive certificate from a supplied dictionary}
Let $f_1,\ldots,f_N$ be probability laws, let $\alpha$ be a probability
row, and assume explicitly
\[
 p=\sum_i\alpha_i f_i.
\]
For nonnegative occupation variables $x_i,z_i,y_{aij}$, impose
\[
 x_j=\alpha_j+\sum_{a,i}y_{aij},\qquad
 x_i=z_i+\sum_{a,j}y_{aij},\qquad \sum_i x_i\le V.
\]
Define the residuals
$R_i=x_if_i-z_i\delta_\varepsilon-\sum_{a,j}y_{aij}L_af_j$.
The local defects must be weighted by how often the decoded process
visits the corresponding states.

\begin{theorem}[Occupation certificate]
For $x_i>0$, decode $T_a(i,j)=y_{aij}/x_i$ and
$\tau_i=z_i/x_i$; use an immediate stop for zero-occupation rows.
The resulting law $q$ started from $\alpha$ is proper and satisfies
\[
 \mathbb E_q|W|\le V-1,\qquad
 \operatorname{TV}(p,q)\le\frac12\sum_i\|R_i\|_1.
\]
If the source mixture is only accurate to $\epsilon_0$ in $\ell^1$,
add $\epsilon_0/2$ to the bound.
\end{theorem}
Theorem~\ref{II:U:thm:flows} proves properness and the telescoping residual
estimate. The budget $V$ controls expected visits; an unweighted maximum
local defect cannot replace the occupation-weighted total in general.

The finite-sample version fixes the dictionary before validation, or uses
an independent split. For each dictionary law, empirical bins up to an
integer horizon $L$ and one overflow bin give a simultaneous confidence
event. On this event a finite LP may optimize its flows using the same
histograms, since the event is uniform over feasible flows. The bound
$\operatorname{TV}(p,\widehat q)\le\widehat C_L/2$ certifies an actually
returned feasible model. This is a high-probability statement under the
sampling model, not an unconditional certificate from arbitrary traces.
The complete oracle inequality is Theorem~\ref{II:U:thm:flow-lp}.

\section{What positive acquisition requires}
The preceding theorem starts with a dictionary and its source-mixture
relation. Discovering a small suitable dictionary is a different problem.
The conditional spectral--conic bridge explicitly combines prefix-mass
bounds, spectral conditioning, a feasible visit budget, and a conic defect.
Its bound charges the whole product of these parameters; a small local
closure defect alone is insufficient.

The constructive results below separate several access models:
a supplied finite candidate pool; positive closure in dual state
coordinates; rank-two monotone endpoint acquisition; observable-coordinate
rounding; relatively separated prefix pools; and bounded-depth promise
classes. Each theorem states its own conditioning, positivity, sampling
and feasibility assumptions. None supplies unrestricted discovery from
uninterpreted observations.

A two-state proper law already prevents replacing latent positive states
by an arbitrary finite set of actual prefix residuals with exact closure.
The golden-ratio example makes this obstruction quantitative. Additional
counterexamples show that uniform Doeblin positivity does not imply
arbitrary hidden-state purification and that changes to a rare-prefix
posterior invalidate naive conditional smoothing bounds. These examples
are part of the theorem scope, rather than failures of an implementation.

\section{Word laws, conditional futures and filtering}
A total-variation guarantee for complete words does not directly become
a uniform next-event guarantee after every prefix. Conditioning divides
by prefix mass. Likewise, a projective contraction statement for a filter
state and a KL statement for the law of an entire word are different
claims. Projective contraction and stability of nonlinear filters have
a substantial prior literature~\citep{cohen,atar}. The final proofs in this section give the sharp projective-to-KL
envelope, finite-horizon word-law estimates, and the precise columnwise
limit hypothesis used for word laws. Intermittent contraction controls the
filter state under its own assumptions.

The open acquisition problem is to find, with quantified sample and bit
costs, a small dictionary satisfying the hypotheses that make the entire
certificate useful. A successful theorem must account for rare-prefix
sampling, adaptive dictionary selection and numerical feasibility. The
transfer-operator constructions in Section~\ref{sec:memory} assume
a supplied dynamical model.

\section{Reading guide}
The proofs below use the linear preliminaries and attainment result in
Subsections~\ref{I:U:sec:preliminaries} and~\ref{I:U:sec:realization}.
They begin with the three KL regimes and the fixed-state obstruction.
The occupation and positive-realization section contains all constructive
and conditional acquisition results. The final sections give projective
bounds and exact counterexamples. Numerical agreement is not used to
establish global optimality or bibliographic priority.
\label{II:II:main-end}

\section*{Detailed proofs and supporting results}

\section{Probability guarantees: rank one and fixed state budgets}\label{II:U:sec:probability}
This section concerns probability distributions on \emph{completed finite
words}. All logarithms are natural. Let $p$ be such a law, $d\ge2$, and
$p(\varepsilon)\ge\mu>0$. Define
\[
 \mathcal Q_1=\left\{q_{a,b}(w)=a\prod_{i=1}^{|w|}b_{w_i}:
 a>0,\ b_x\ge0,\ a+\sum_xb_x=1\right\},\quad
 \mathcal K(p)=\inf_{q\in\mathcal Q_1}D_{\rm KL}(p\|q).
\]
An impulse word-Hankel is a finite partial permutation of norm one.
Consequently $H_p$ is a norm-convergent sum of finite-rank operators and is
compact with $\|H_p\|\le1$. A nonzero normalized nonnegative rank-one
word-Hankel has a positive empty pivot and multiplicative coefficients;
its law is in $\mathcal Q_1$.

\subsection{The information projection is explicit}
Let $M=\mathbb E_p|W|$, $c_x=\mathbb E_p\#_x(W)$, and
$g(M)=(1+M)\log(1+M)-M\log M$, with $g(0)=0$.
\begin{proposition}[Classical one-state information projection]\label{II:U:prop:iprojection}
If $M<\infty$, the optimum is attained at
\[
 a_*=(1+M)^{-1},\qquad b_{*,x}=c_x/(1+M),
\]
and, for $M>0$ and $\pi_x=c_x/M$,
\[
 \mathcal K(p)=g(M)+M H(\pi)-H(p).
\]
The value is zero when $M=0$.
\end{proposition}
\begin{proof}
Cross-entropy is $-\log a-\sum_xc_x\log b_x$ on the event simplex.
The entropy maximization or Lagrange multiplier identity normalizes the
counts $(1,c_x)$ to obtain the displayed parameters. A zero $c_x$ means
that letter is absent almost surely, so it causes no support problem.
Finite mean implies $H(p)\le g(M)+M\log d<\infty$: first maximize
entropy at fixed length, then use the geometric maximum-entropy law at
fixed mean length. Subtracting $H(p)$ gives the formula.
\end{proof}

This proposition is an ingredient, not a novelty claim. The substantive
question is how small $\sigma_2(H_p)$ forces $\mathcal K(p)$ to be.

\subsection{From a spectral tail to a proper approximation}
Write $t_p(T)=\Pr_p\{|W|>T\}$ and $N_d(T)=\sum_{j=0}^Td^j$.
\begin{lemma}[Spectral-to-TV bridge]\label{II:U:lem:spectral-tv}
If $\epsilon=\sigma_2(H_p)\le\mu/2$, there is a proper
$q\in\mathcal Q_1$ such that, for every integer $T\ge0$,
\[
 \UTV(p,q)\le \frac{dC_\mu}{\mu}\epsilon\sqrt{N_d(T)}
                    +\frac{t_p(T)}{\mu^2},\qquad
 C_\mu=(1+1/\mu)(1+2/\mu).
\]
It can be chosen with stop probability $p(\varepsilon)$, hence with
$\Pr_q\{|W|>T\}\le(1-\mu)^{T+1}$.
\end{lemma}
\begin{proof}
Set $a=p(\varepsilon)$ and $c_x=p(x)/a$. Take a leading rank-one SVD
term $H_p=su v^*+E$ with $\|E\|=\epsilon$ and $s\le1$.
The empty pivot implies $|v_\varepsilon|\ge a-\epsilon\ge\mu/2$.
For $z_x=e_x-c_xe_\varepsilon$, the row
$h_x(v')=p(xv')-c_xp(v')$ satisfies $z_x^*H_pe_\varepsilon=0$.
Eliminating the rank-one term using this equality gives
$\|h_x\|_2\le\epsilon\|z_x\|(1+|v_\varepsilon|^{-1})
\le C_\mu\epsilon$.

On words of length at most $T$, sum Cauchy--Schwarz over the $d$ rows.
On longer words, positivity gives
\[
 \sum_x\sum_{|v'|>T}|h_x(v')|
 \le t_p(T+1)+\left(\sum_xc_x\right)t_p(T)
 \le t_p(T)/a.
\]
The last inequality uses $\sum_xp(x)\le1-a$. Put
$b_x=\sum_{v'}p(xv')$, so $\sum b_x=1-a$ and
$b_x-c_x=\sum h_x$. The defect in
$p=a\delta_\varepsilon+Kp$, where $(Kf)(xv')=b_xf(v')$, has
$\ell^1$ norm at most $2\sum_x\|h_x\|_1$.
The proper law $q=a\prod b_{w_i}$ solves that recurrence exactly.
Since $\|K\|_{1\to1}=1-a$, its inverse resolvent has norm at most
$1/a$. Divide the resulting $\ell^1$ bound by two and use $a\ge\mu$.
\end{proof}


\subsection{Three sharp regimes}
For fixed $C>0$, $0<\rho<1$, and $0<\mu<1$, let
\[
 \mathcal P_{\exp}=\{p:p(\varepsilon)\ge\mu,
       \ t_p(t)\le C\rho^{t+1}\text{ for every integer }t\ge0\},
 \quad \gamma=\frac{\log(1/\rho)}{\log(1/\rho)+\tfrac12\log d}.
\]
\begin{theorem}[Sharp exponential-tail rate]\label{II:U:thm:kl-exp}
As $\epsilon\downarrow0$,
\[
 \sup_{\substack{p\in\mathcal P_{\exp}\\\sigma_2(H_p)\le\epsilon}}
 \mathcal K(p)
 =\Theta\!\left(\epsilon^\gamma\log^2(1/\epsilon)\right).
\]
The asymptotic range may be chosen with $0<\epsilon\le\mu/2$ and with the TV upper bound $0<\Delta\le1/2$ used below; all truncation horizons are integers. The constants depend on $d,\mu,C,\rho$. The claim is for the best
normalized one-state law, not a particular spectral estimator.
\end{theorem}
\begin{proof}[Upper bound]
Restrict to sufficiently small $0<\epsilon\le\mu/2$ and put
$\lambda=\log(1/\rho)$ and
$T=\lfloor\log(1/\epsilon)/(\lambda+\tfrac12\log d)\rfloor\in\mathbb Z_{\ge0}$ in
Lemma~\ref{II:U:lem:spectral-tv}. Both terms are $O(\epsilon^\gamma)$.
Let $\Delta\le1/2$ be a positive such upper bound. Smooth the event
law $\pi=(a,b_1,\ldots,b_d)$ to
$\pi_\Delta=(1-\Delta)\pi+\Delta\operatorname{Unif}(d+1)$,
and call the resulting proper law $q_\Delta$.

Couple $W\sim p,V\sim q$ maximally, enlarging the mismatch event $E$
to probability $\Delta$ if necessary. Exponential tails imply
\[
 \mathbb E[|W|1_E]\le\Delta A_p(\Delta),\quad
 \mathbb E[|V|1_E]\le\Delta A_q(\Delta),\quad A_p,A_q=O(\log(1/\Delta)).
\]
For example, summing $\min(\Delta,C\rho^{t+1})$ proves the first
bound. With $A=\max(A_p,A_q)$, conditional entropy and the geometric
maximum-entropy bound imply
\[
 |H(p)-H(q)|\le h_2(\Delta)+\Delta\{g(A)+A\log d\}.
\]
All smoothed event probabilities are at least $\Delta/(d+1)$. The stop
term cancels when comparing two word cross-entropies, giving a difference
at most $\Delta(A_p+A_q)\log((d+1)/\Delta)$.
Finally the event-by-event renewal calculation yields
\[
 D_{\rm KL}(q\|q_\Delta)=a^{-1}D_{\rm KL}(\pi\|\pi_\Delta)
 \le-\mu^{-1}\log(1-\Delta).
\]
Adding the three estimates gives $D_{\rm KL}(p\|q_\Delta)
=O(\Delta\log^2(1/\Delta))$. Optimization can only reduce it.
\end{proof}

\begin{proof}[Lower bound]
For an integer $T\ge2$ take
$p_T=(1-\delta)\delta_\varepsilon+\delta\operatorname{Unif}(\Sigma^T)$.
In the orthonormal uniform shell basis of lengths $0,\ldots,T$, its
Hankel is $bJ+(1-\delta)e_0e_0^*$, where $J$ reverses the coordinates
and $b=\delta d^{-T/2}$. All orthogonal shell directions vanish.
The endpoint block is
$\left(\begin{smallmatrix}1-\delta&b\\b&0\end{smallmatrix}\right)$;
one singular value exceeds $b$, one is smaller, and the $T-1$ interior
singular values equal $b$. Hence $\sigma_2=b$ exactly.
Proposition~\ref{II:U:prop:iprojection} gives
\[
 \mathcal K(p_T)=g(\delta T)-h_2(\delta).
\]
Choose $\delta=c_0\rho^T$, $c_0=\min(C,1)$. The tail constraint holds,
and $p_T(\varepsilon)\ge\mu$ for large $T$.
Expanding at zero gives
$g(\delta T)-h_2(\delta)\sim\lambda c_0\rho^TT^2$.
Since $\epsilon_T=c_0(\rho/\sqrt d)^T$, this is the required order.
Successive $\epsilon_T$ have a fixed ratio, so selecting the first one
below a given tolerance extends the bound from a subsequence to all small
$\epsilon$.
\end{proof}


\begin{theorem}[Only a mean bound: exact discontinuity]\label{II:U:thm:kl-mean}
Fix $M_0>0$ and $0<\mu<1$. For every $\epsilon>0$,
\[
 \sup_{\substack{p(\varepsilon)\ge\mu,\ \mathbb E|W|\le M_0\\
                      \sigma_2(H_p)\le\epsilon}}\mathcal K(p)
 =g(M_0)+M_0\log d.
\]
At $\epsilon=0$ this supremum is zero.
\end{theorem}
\begin{proof}
The information-projection formula, $H(p)\ge0$, $H(\pi)\le\log d$,
and monotonicity of $g$ give the upper bound. Let $T$ run through
positive integer multiples of $d$ with $T\ge M_0/(1-\mu)$, let $w_T$ have equal symbol counts, and put
$p_T=(1-M_0/T)\delta_\varepsilon+(M_0/T)\delta_{w_T}$.
Its mean is $M_0$, its entropy is $h_2(M_0/T)$, and its expected-count
law is uniform. Removing the rank-one empty impulse shows
$\sigma_2(H_{p_T})\le M_0/T$, since an impulse Hankel has norm one.
Thus its optimal KL tends to the upper bound while satisfying any
positive spectral tolerance. If $\sigma_2=0$, positivity of the empty
pivot makes $p$ itself a proper multiplicative rank-one law.
\end{proof}


\begin{theorem}[Higher length moments]\label{II:U:thm:kl-moment}
Fix $s>1$, $M_s>0$, and $0<\mu<1$. Then
\[
 \sup_{\substack{p(\varepsilon)\ge\mu,\ \mathbb E|W|^s\le M_s\\
                      \sigma_2(H_p)\le\epsilon}}\mathcal K(p)
 =\Theta\!\left((\log(1/\epsilon))^{1-s}
                     \log\log(1/\epsilon)\right).
\]
This is an $\epsilon\downarrow0$ assertion; take $0<\epsilon\le\mu/2$ sufficiently small that the TV bound $\Delta$ in the proof is at most $1/2$. All truncation and shell lengths are integers. The constants are not uniform as $s\downarrow1$.
\end{theorem}
\begin{proof}
Markov's inequality and Lemma~\ref{II:U:lem:spectral-tv}, with
$T=\lfloor\log(1/\epsilon)/\log d\rfloor$, give a proper $q$ with
$\UTV(p,q)\le\Delta=O((\log(1/\epsilon))^{-s})$.
Its length is geometric with stop at least $\mu$, so its $s$th moment
is uniformly bounded. In the preceding coupling, H\"older gives
$\mathbb E[(|W|+|V|)1_E]=O(\Delta^{1-1/s})$.
The same entropy and smoothed cross-entropy argument gives KL
$O(\Delta^{1-1/s}\log(1/\Delta))$.

For the lower bound use the uniform shell source with $\delta=M_s/T^s$
and large integer $T$ with $M_s/T^s\le1-\mu$. Its $s$th moment is exactly $M_s$,
$\epsilon_T=M_sT^{-s}d^{-T/2}$, and
\[
 \mathcal K(p_T)=g(M_sT^{1-s})-h_2(M_sT^{-s})
 \sim M_s(s-1)T^{1-s}\log T.
\]
Here $\log(1/\epsilon_T)\sim T\log d/2$; consecutive spectral values
have ratios bounded away from zero. The asserted order follows.
\end{proof}


\subsection{Why the probability theorem does not contradict the obstruction}
The binary family in Subsection~\ref{I:U:sec:obstruction} is not the fixed-tail
probability class above. Here rank-one positivity and the nonvanishing empty
pivot give a special closure recurrence; the tail condition converts an
$\ell^2$ defect to $\ell^1$ and controls entropy. None of those steps gives a
general rank-$k$ proper realization with $k$ positive states. Nor does a
small completed-word TV guarantee uniform next-token error after every
prefix: conditioning divides by the prefix probability. Those are separate
positive-realization and conditional-prediction questions.
\subsection{A KL discontinuity at every fixed positive-state budget}
\label{II:U:sec:kl-fixed-states}
Let $\mathcal Q_N^+$ be the laws generated by proper probabilistic
automata with at most $N$ states. Concretely,
$q(w)=\alpha T_w\tau$, where $\alpha\ge0$, $\alpha\mathbf1=1$,
$T_a\ge0$, $\tau\ge0$, and
$\tau+\sum_aT_a\mathbf1=\mathbf1$; termination is almost sure from
$\alpha$. Zero transition weights and unused states are permitted.
Write $\mathcal K_N(p)=\inf_{q\in\mathcal Q_N^+}D_{\rm KL}(p\|q)$,
so that $\mathcal K_1=\mathcal K$.

\begin{theorem}[Mean-constrained KL discontinuity for fixed $N\ge2$]
\label{II:U:thm:kl-fixed-states}
Fix $d\ge2$, $M_0>0$, and $\mu\in(0,1)$. For an integer $m\ge1$,
put $K=d^m$ and $L=mK$, and concatenate the $K$ distinct words of
$\Sigma^m$, each exactly once and in any order, to form $w_L$.
Whenever $L>M_0/(1-\mu)$, define
\[
 \delta_L=\frac{M_0}{L},\qquad
 p_L=(1-\delta_L)\delta_\varepsilon+\delta_L\delta_{w_L}.
\]
Then $\mathbb E_{p_L}|W|=M_0$, $p_L(\varepsilon)\ge\mu$, and
$\sigma_{k+1}(H_{p_L})\le\sigma_2(H_{p_L})\le M_0/L$ for every
$k\ge1$. For every integer $N\ge2$,
\begin{align}
 M_0\log d-\frac{M_0\log N}{m}-h_2(\delta_L)
 &\le\mathcal K_N(p_L)\notag\\
 &\le M_0\log d+\frac{M_0}{L}
 \left[\log L-(L-1)\log(1-1/L)\right],
 \label{II:U:eq:kl-fixed-sandwich}
\end{align}
where $h_2(s)=-s\log s-(1-s)\log(1-s)$. Consequently, for each
\emph{fixed} $N\ge2$,
\[
 \lim_{m\to\infty}\mathcal K_N(p_L)=M_0\log d>0.
\]
For the same family, $\mathcal K_1(p_L)\to g(M_0)+M_0\log d$,
whereas $L+1$ states realize $p_L$ exactly.
\end{theorem}
\begin{proof}
The mass and mean assertions are immediate. The impulse Hankel
$H_{\delta_{w_L}}$ is a partial permutation of norm one. Subtracting
the rank-one operator $(1-\delta_L)H_{\delta_\varepsilon}$ therefore
gives $\sigma_2(H_{p_L})\le\delta_L$.

For any $q\in\mathcal Q_N^+$, pad its realization with unused states
if necessary, and let $T=\sum_aT_a$. Since $T\mathbf1\le\mathbf1$,
also $T^m\mathbf1\le\mathbf1$. For $B\in\Sigma^m$ put
$\psi(B)=\|T_B\|_{\infty\to\infty}
=\max_i(T_B\mathbf1)_i$. Nonnegativity yields
\[
 \sum_{B\in\Sigma^m}\psi(B)
 \le\sum_i\sum_{B\in\Sigma^m}(T_B\mathbf1)_i
 =\sum_i(T^m\mathbf1)_i\le N.
\]
If $w_L=B_1\cdots B_K$, submultiplicativity, $\|\alpha\|_1=1$,
and $\|\tau\|_\infty\le1$ imply
\[
 q(w_L)\le\prod_{r=1}^K\psi(B_r)
 \le\left(\frac NK\right)^K.
\]
The last inequality is arithmetic--geometric mean; it remains valid
when a factor vanishes. Thus
$-\log q(w_L)\ge L\log d-(L/m)\log N$.
Using $q(\varepsilon)\le1$ gives
\[
 D_{\rm KL}(p_L\|q)
 \ge-h_2(\delta_L)-\delta_L\log q(w_L),
\]
which proves the lower bound, including the case of infinite KL.

For the upper bound use the following two-state automaton:
\[
 \alpha=(1,0),\qquad
 \tau=\begin{pmatrix}1-M_0/L\\1/L\end{pmatrix},\qquad
 T_a=\begin{pmatrix}
 0&M_0/(Ld)\\0&(1-1/L)/d
 \end{pmatrix}\quad(a\in\Sigma).
\]
It is proper, has mean length $(M_0/L)L=M_0$, and satisfies
\[
 q_L(\varepsilon)=1-M_0/L,\qquad
 q_L(w)=\frac{M_0}{L^2}(1-1/L)^{L-1}d^{-L}
 \quad(|w|=L).
\]
Its divergence is exactly the upper expression in
\eqref{II:U:eq:kl-fixed-sandwich}. The two residual terms in that expression
vanish, as do $m^{-1}\log N$ and $h_2(\delta_L)$ for fixed $N$.

Each letter occurs $md^{m-1}=L/d$ times in $w_L$. Hence
Proposition~\ref{II:U:prop:iprojection} gives
$\mathcal K_1(p_L)=g(M_0)+M_0\log d-h_2(\delta_L)$.
Finally, a deterministic path of length $L$, with initial stop
probability $1-\delta_L$, initial continuation probability $\delta_L$,
and terminal stop probability one, realizes $p_L$ with $L+1$ states.
\end{proof}


\paragraph{Quantifiers and comparison classes.}
For every fixed $N\ge2$ and every positive spectral tolerance, the
mean-constrained worst-case value of $\mathcal K_N$ is at least
$M_0\log d$. At tolerance $\sigma_2(H_p)=0$, it is zero because
$p\in\mathcal Q_1\subset\mathcal Q_N^+$. The theorem identifies the
exact limiting value on the displayed family, not the exact worst-case
supremum over all sources at a positive tolerance. It also does not assert
that every source with small spectral tail has large KL error. The
quantifier that $N$ is fixed before $m\to\infty$ is essential.
Replacing $\delta_{w_L}$ by the uniform law on $\Sigma^L$ changes the
example: the same two-state $q_L$ then has divergence
$\frac{M_0}{L}[\log L-(L-1)\log(1-1/L)]\to0$.

\section{Positive realization through occupation flows}\label{II:U:sec:flows}
The second probability bridge starts from a finite dictionary rather than a
singular value. Let $f_1,\ldots,f_N$ be probability laws on finite words,
$\alpha$ a probability row vector, and initially assume
$p=\sum_i\alpha_i f_i$. Let $(L_af)(av)=f(v)$ and zero elsewhere.
For a visit budget $V\ge1$, choose nonnegative $x_i,z_i,y_{aij}$ with
\begin{equation}\label{II:U:eq:flow}
 x_j=\alpha_j+\sum_{a,i}y_{aij},\qquad
 x_i=z_i+\sum_{a,j}y_{aij},\qquad \sum_i x_i\le V.
\end{equation}
These imply $\sum z_i=1$ and $x\ge\alpha$. Define
\[
 R_i=x_if_i-z_i\delta_\varepsilon-\sum_{a,j}y_{aij}L_af_j,
 \qquad J_\infty=\sum_i\|R_i\|_1.
\]

\begin{theorem}[Properness and the occupation-weighted certificate]\label{II:U:thm:flows}
For $x_i>0$ set $T_a(i,j)=y_{aij}/x_i$ and $\tau_i=z_i/x_i$.
For $x_i=0$, take a zero transition row and $\tau_i=1$.
The automaton started at $\alpha$ terminates almost surely and defines
a normalized law $q$ with
\[
 \mathbb E_q|W|\le V-1,\qquad \UTV(p,q)\le J_\infty/2.
\]
If $\|p-\sum_i\alpha_if_i\|_1\le\epsilon_0$, add
$\epsilon_0/2$ to the TV bound.
\end{theorem}
\begin{proof}
Let $T=\sum_aT_a$. The constraints give $\tau+T\mathbf1=\mathbf1$
and $x=\alpha+xT$. Iteration shows
$\nu=\sum_{t\ge0}\alpha T^t\le x$. Survival probabilities
$\alpha T^t\mathbf1$ have finite sum, hence tend to zero. Thus
termination holds from $\alpha$ and
$1+\mathbb E|W|=\nu\mathbf1\le V$.

For $x_i>0$ write $r_i=R_i/x_i$, and let $q_i$ be the possibly
subprobability law starting in state $i$. The two recurrences imply
\[
 f_i-q_i=r_i+\sum_{a,j}T_a(i,j)L_a(f_j-q_j).
\]
After $t$ expansions and weighting by $\alpha$, the remainder has norm
at most $2\alpha T^t\mathbf1\to0$. Prepending is an $\ell^1$
isometry, so
$\|p-q\|_1\le\sum_i\nu_i\|r_i\|_1\le J_\infty$.
A zero-visit state has no incoming flow and $\alpha_i=0$, hence makes
no contribution. Unreachable circulation in $x-\nu$ is harmless.
The misspecification term follows by the triangle inequality.
\end{proof}


\paragraph{Sharpness of the visit factor.}
Take a state looping on $c$ with probability $\beta<1$, exiting on $a,b$
to an immediate-stop state with probabilities $\gamma$ each, and otherwise
stopping; let $0<2\gamma\le1-\beta$. Change the two source exit
probabilities to $\gamma+\eta,\gamma-\eta$, $0<\eta\le\gamma$.
Using the source future law and the terminal law as dictionary gives
$r_1=\eta(\delta_a-\delta_b)$, $r_2=0$, and
$\nu_1=(1-\beta)^{-1}$. The discrepancies occur on disjoint pairs
$c^na,c^nb$, so
$\UTV(p,q)=\eta/(1-\beta)=\nu_1\|r_1\|_1/2$.
The factor cannot be replaced by a uniform unweighted local error.

\subsection{A finite-sample linear program}\label{II:U:sec:flow-lp}
Fix $L\ge1$. For each dictionary law take $n_i$ independent samples and
record words up to length $L$ plus an overflow category. The dictionary
must be fixed before these samples, or chosen on an independent split.
Independence between the sample sets for different dictionary entries is
not needed. Let
\[
 K_L=1+\sum_{j=0}^Ld^j,\qquad
 b_i=\min\left\{2,\sqrt{\frac{2(K_L\log2+\log(N/\delta))}{n_i}}\right\}.
\]
A union bound over all sign vectors and all dictionary entries, with
Hoeffding's inequality, gives simultaneous histogram $\ell^1$ errors
at most $b_i$ with probability at least $1-\delta$.

Use hats for empirical probabilities and tails. For feasible flows put
\begin{align*}
 \widehat H_L&=\sum_i\sum_{|w|\le L}
  \left|x_i\widehat f_i(w)-z_i1_{w=\varepsilon}
             -\sum_{a,j}y_{aij}(L_a\widehat f_j)(w)\right|,\\
 \widehat U_L&=\sum_i x_i\widehat t_i(L)
                 +\sum_{a,i,j}y_{aij}\widehat t_j(L-1),\\
 S(x)&=\sum_i x_ib_i+\sum_{a,i,j}y_{aij}b_j
       =\sum_i(2x_i-\alpha_i)b_i,\\
 \widehat C_L&=\widehat H_L+\widehat U_L+S(x).
\end{align*}
Minimizing $\widehat C_L$ subject to \eqref{II:U:eq:flow} is a finite LP,
using the usual epigraph variables for absolute values.

\begin{theorem}[Uniform data-dependent certificate and defect oracle]\label{II:U:thm:flow-lp}
Assume $f_i$ are probability laws, $\alpha$ is a probability vector and $p=\sum_i\alpha_i f_i$. Fix the dictionary independently of the validation data and use the sampling scheme above. On one event of probability at least $1-\delta$, every feasible decoded
law satisfies $\UTV(p,q)\le\widehat C_L/2$. In particular the LP
solution $\widehat q$ satisfies
\[
 \UTV(p,\widehat q)\le
 \inf_{\text{feasible flows}}\left\{\tfrac12J_\infty
                     +\sum_i x_it_i(L)+S(x)\right\}.
\]
If $t_i(L)\le t_{\max}(L)$ and $b_i\le b$, this is at most
$\tfrac12\inf J_\infty+Vt_{\max}(L)+(2V-1)b$.
\end{theorem}
\begin{proof}
Define $J_L=H_L+U_L$ from the true probabilities. Tail majorization
gives $J_\infty\le J_L$. Head bins and the overflow bin are disjoint;
each shifted suffix histogram is a coarsening of the same sampled
histogram. Consequently
$|J_L-\widehat H_L-\widehat U_L|\le S(x)$ for \emph{all} feasible
flows on the simultaneous event. Theorem~\ref{II:U:thm:flows} proves the
first assertion, including for fitted flows chosen from these data.

For a comparison flow, optimality implies
$\widehat C_L^*\le J_L+2S(x)$. For nonnegative tail vectors $u,v$,
$\|u\|_1+\|v\|_1-\|u-v\|_1\le2\|u\|_1$.
Applying this with the source tail gives
$J_L\le J_\infty+2\sum_i x_it_i(L)$. Divide by two and infimize.
Finally $\sum x_i\le V$ and $\sum\alpha_i=1$ bound $S(x)$.
\end{proof}


The oracle compares occupation-weighted \emph{closure defects in the supplied
dictionary}, not the best arbitrary $N$-state TV or KL approximation.
Cancellations can make actual error smaller than the certificate. An
additive LP objective gap $\zeta$ costs at most $\zeta/2$ in TV if the
returned flow is feasible. Numerical feasibility itself must be checked or
repaired. The LP has size $O(dN^2+NK_L)$ up to constants; $K_L$ is
exponential in $L$. Under uniform exponential dictionary tails and
$n_i\ge n$, balancing terms gives a rate of order
\[
 Vn^{-\lambda/(2\lambda+\log d)}
       +V\sqrt{\log(N/\delta)/n},\qquad\lambda=\log(1/\rho),
\]
apart from the approximation defect and fixed constants. This is not a
dimension-free or minimax-optimal claim.

\subsection{A safe KL extension with an explicitly charged extra state}\label{II:U:sec:flow-smoothing}
For any proper $q$ with $\UTV(p,q)\le\Delta<1$, set
$u(w)=2^{-|w|-1}d^{-|w|}$ and
$q^+=(1-\Delta)q+\Delta u$. This mixture can be realized by adding
one probabilistic state and modifying the initial mixture. Define
$\Psi_p(\Delta)=\sum_{t\ge0}\min(\Delta,t_p(t))$.
\begin{proposition}[Smoothing a completed-word TV approximation]
For finite $\Psi_p(\Delta)$,
\[
 D_{\rm KL}(p\|q^+)\le-\log(1-\Delta)+\Delta\log2
                            +\log(2d)\Psi_p(\Delta).
\]
\end{proposition}
\begin{proof}
Split a common submeasure from $p$ and $q$, of mass $1-\Delta$,
enlarging the excess if their actual TV is smaller. Let the remaining
submeasure of $p$ be $p_1$ of mass $\Delta$ and the common one $p_0$.
The denominator $q^+$ dominates $(1-\Delta)p_0+\Delta u$.
The log-sum inequality bounds the common part by
$-(1-\Delta)\log(1-\Delta)$ and the excess by
$\Delta D_{\rm KL}(p_1/\Delta\|u)$.
Dropping the nonnegative entropy of $p_1/\Delta$ bounds the latter by
$\Delta\log2+\log(2d)\sum_w|w|p_1(w)$.
The last sum is at most $\Psi_p(\Delta)$ by summing tail probabilities.
Finally replace the common term by $-\log(1-\Delta)$.
\end{proof}

For exponential source tails this is $O(\Delta\log(1/\Delta))$.
The extra state changes the approximation class; it does not improve the
sharp one-state rate in Theorem~\ref{II:U:thm:kl-exp} within the same class.

\paragraph{Keeping the same positive-state count.}
Alternatively, for $0<\Delta<1$ smooth each state locally:
\[
 \tau_i^{(\Delta)}=(1-\Delta)\tau_i+\frac\Delta{d+1},\qquad
 T_a^{(\Delta)}=(1-\Delta)T_a+\frac\Delta{d+1}I_N.
\]
The rows remain stochastic and every state has stop probability at least
$\Delta/(d+1)$. The all-original and all-smoothing choices in the expanded
product give the sum of two nonnegative contributions:
\[
 q_\Delta(w)\ge(1-\Delta)^{|w|+1}q(w)
                   +\left(\frac\Delta{d+1}\right)^{|w|+1}.
\]
Applying log-sum to the common and excess submeasures as before gives,
when $\UTV(p,q)\le\Delta$,
\begin{align*}
 D_{\rm KL}(p\|q_\Delta)
 &\le-\log(1-\Delta)(1+\mathbb E_p|W|)\\
 &\quad+\Delta\log(d+1)+\Psi_p(\Delta)\log\frac{d+1}{\Delta}.
\end{align*}
For exponential tails this is $O(\Delta\log^2(1/\Delta))$.
Sharpness for a general fixed $N$ is not established by this construction.
Local smoothing can increase the original visit budget and changes the
transition topology; both effects must be accounted for.

\subsection{The unresolved acquisition problem}\label{II:U:sec:acquisition-obstruction}
For a prefix anchor $u_i$, a residual dictionary law is
$f_i(v)=p(u_iv)/P_p(u_i\Sigma^*)$. The denominator is a prefix mass,
not the completed-word value $p(u_i)$. Rare anchors may require many
samples; source tails do not provide uniform residual-tail constants
without accounting for that denominator. The supplied-dictionary theorem
does not discover a small nearly closed dictionary or bound its acquisition
cost. Linear rank and positive state size also differ: nonnegative rational
series need not have finite positive realizations. A rank bound alone cannot
replace the positive dictionary hypothesis.

\begin{proposition}[A classical positive-realization obstruction]\label{II:U:prop:positive-obstruction}
For $0<x<1$ and binary words, put $d(w)=\#_a(w)-\#_b(w)$ and
\[
 p_x(w)=(1-x)^2(x/2)^{|w|}(1+d(w))^2.
\]
This is a probability law of counting-Hankel rank exactly three, with
positive empty mass and an exponential tail, but no finite nonnegative
weighted-automaton realization.
\end{proposition}
\begin{proof}
Over all words of length $t$, the mean of $d(w)$ is zero and its second
moment is $t$. Therefore the length mass is
$(1-x)^2(t+1)x^t$, summing to one. The tail has an exponential bound
with any fixed rate strictly between $x$ and one.
Since $d(uv)=d(u)+d(v)$, expanding the square gives rank at most three.
Restriction to words $a^j$ gives a nondegenerate quadratic times a
geometric sequence and hence rank at least three. Its zero language is
$\{w:\#_b(w)=\#_a(w)+1\}$, which is nonregular: intersect with
$a^*b^*$ to obtain $\{a^nb^{n+1}:n\ge0\}$.
A finite automaton with nonnegative weights has regular nonzero support,
because a coefficient is positive exactly when a positive-weight path
exists. Its zero language is consequently regular as well, a contradiction.
\end{proof}

This is a classical type of signed-versus-positive realization obstruction
\citep{denis2008}; this is not a new impossibility claim.
It already has $\sigma_4(H_{p_x})=0$ and shows why the higher-rank bridge
cannot rely on a spectral tail alone to guarantee an exact finite positive
realization.
\subsection{Conditioned residual anchors and signed spectral closure}
\label{II:U:sec:anchored-closure}
The following bridge is conditional on a chosen dictionary and a positive
flow defect. It supplies the previously missing explicit hypotheses; it
does not derive positivity or a small dictionary from singular values.
All data in this subsection are real, and $p$ is a probability law.

\begin{definition}[Conditioned residual anchors]
\label{II:U:def:conditioned-anchors}
Fix $k\ge1$, $N\ge k$, and constants $\mu_{\min}>0$ and
$\kappa<\infty$. Choose prefix anchors
$\mathcal U=(u_1,\ldots,u_N)$ with $u_1=\varepsilon$ and
\[
 \mu_i=P_p(u_i\Sigma^*)=\sum_v p(u_iv)\ge\mu_{\min},\qquad
 f_i(v)=\frac{p(u_iv)}{\mu_i}.
\]
Each $f_i$ is a probability law, $f_1=p$, and, with
$E_{\mathcal U}=[e_{u_1},\ldots,e_{u_N}]$ and
$D_\mu=\Udiag(\mu_1,\ldots,\mu_N)$,
\[
 F_{\mathcal U}=[f_1,\ldots,f_N]
 =H_p^*E_{\mathcal U}D_\mu^{-1}.
\]
Write $H_p=U_k\Sigma_kV_k^*+E_k$ for a truncated SVD and put
$\epsilon_k=\|E_k\|=\sigma_{k+1}(H_p)$. The spectral coordinate
matrix must satisfy
\begin{equation}
 C_{\mathcal U}=V_k^*F_{\mathcal U}
 =\Sigma_kU_k^*E_{\mathcal U}D_\mu^{-1},\qquad
 \Urank C_{\mathcal U}=k,\qquad
 \|C_{\mathcal U}^{\dagger}\|_{2\to2}
 =\frac1{\sigma_k(C_{\mathcal U})}\le\kappa.
 \label{II:U:eq:anchor-coordinates}
\end{equation}
These conditions define
$\mathcal H_{\rm anchor}(k,N,\mu_{\min},\kappa)$.
\end{definition}

The orientation is forced by \eqref{I:U:eq:intertwining}: its adjoint is
$L_a^*H_p^*=H_p^*R_a$. Therefore
\begin{equation}
 f_i=\mu_i^{-1}H_p^*e_{u_i},\qquad
 L_a^*f_i=\mu_i^{-1}H_p^*R_ae_{u_i}
 =\mu_i^{-1}H_p^*e_{u_i a},\qquad
 (L_a^*f_i)(v)=f_i(av).
 \label{II:U:eq:anchor-orientation}
\end{equation}
The future subspace is the right singular subspace of $H_p$; using its
left singular vectors here would interchange the prefix and suffix roles.

\begin{lemma}[Conditioned signed linear closure]
\label{II:U:lem:signed-anchor-closure}
Under Definition~\ref{II:U:def:conditioned-anchors}, set
\[
 w_{a,i}^{\rm lin}=C_{\mathcal U}^{\dagger}V_k^*(L_a^*f_i).
\]
These coefficients, which need not be nonnegative, satisfy
\begin{equation}
 \|w_{a,i}^{\rm lin}\|_2\le\kappa,\qquad
 \left\|L_a^*f_i-F_{\mathcal U}w_{a,i}^{\rm lin}\right\|_2
 \le c_{\mathcal U}\epsilon_k,\qquad
 c_{\mathcal U}=\frac{1+\kappa\sqrt N}{\mu_{\min}}.
 \label{II:U:eq:signed-anchor-defect}
\end{equation}
This is linear closure, not affine or positive closure.
\end{lemma}
\begin{proof}
Put $P=V_kV_k^*$ and $g=L_a^*f_i$. Full row rank of
$C_{\mathcal U}$ gives $PF_{\mathcal U}w_{a,i}^{\rm lin}=Pg$, so
\[
 g-F_{\mathcal U}w_{a,i}^{\rm lin}
 =(I-P)g-(I-P)F_{\mathcal U}w_{a,i}^{\rm lin}.
\]
By \eqref{II:U:eq:anchor-orientation}, $\|(I-P)g\|_2\le
\epsilon_k/\mu_{\min}$. Each column of $(I-P)F_{\mathcal U}$
has the same bound, hence its operator norm is at most
$\sqrt N\epsilon_k/\mu_{\min}$.
Finally $\|g\|_2\le\|f_i\|_1=1$ yields
$\|w_{a,i}^{\rm lin}\|_2\le\kappa$. The triangle inequality proves
\eqref{II:U:eq:signed-anchor-defect}.
\end{proof}


\subsection{The conic hypothesis and the conditional higher-rank bridge}
\label{II:U:sec:conditional-cone-bridge}
Fix the anchored dictionary above and $\alpha=(1,0,\ldots,0)$.
For $V\ge1$, let $\mathcal F_N(V)$ consist of nonnegative flows
$(x,z,y)$ satisfying \eqref{II:U:eq:flow}; explicitly,
\[
 x_j=\alpha_j+\sum_{a,i}y_{aij},\qquad
 x_i=z_i+\sum_{a,j}y_{aij},\qquad \sum_i x_i\le V.
\]
This is a nonempty compact polytope. In particular,
$\sum_i z_i=1$ and $\sum_{a,i,j}y_{aij}=\sum_i x_i-1$.

\begin{definition}[Spectral-reference conic defect in $\ell^2$]
\label{II:U:def:cone-defect}
For a feasible flow, define
\begin{align}
 \eta_2(x,z,y)
 &=\sum_i|x_if_i(\varepsilon)-z_i|\notag\\
 &\quad+\sum_{a,i}
 \left\|F_{\mathcal U}
       (x_iw_{a,i}^{\rm lin}-y_{a,i,\cdot})\right\|_2,
 \label{II:U:eq:cone-defect}
\end{align}
where $y_{a,i,\cdot}=(y_{ai1},\ldots,y_{aiN})^\top$.
The condition $\mathcal H_{\rm cone}^{(2)}(V,\eta_{\rm cone})$
means that some flow in $\mathcal F_N(V)$ has
$\eta_2(x,z,y)\le\eta_{\rm cone}$.
Equivalently, the attained minimum
$\eta_2^*(\mathcal U,V)=\min_{\mathcal F_N(V)}\eta_2$ is at most
$\eta_{\rm cone}$. This hypothesis is separate from
$\mathcal H_{\rm anchor}$ and the spectral bound.
\end{definition}

\paragraph{SOCP versus LP.}
Let $G=F_{\mathcal U}^*F_{\mathcal U}$, so
$G_{jj'}=\langle f_j,f_{j'}\rangle_{\ell^2}$ and
$\|F_{\mathcal U}c\|_2=\|G^{1/2}c\|_2$.
With $G$, $w_{a,i}^{\rm lin}$, and $f_i(\varepsilon)$ fixed and
known, epigraph constraints
$\|G^{1/2}(x_iw_{a,i}^{\rm lin}-y_{a,i,\cdot})\|_2\le s_{ai}$,
together with linear epigraphs for the absolute values, give a finite
second-order cone program (SOCP). Singularity of $G$ is harmless.
An arbitrary infinite Hankel oracle does not automatically supply these
finite coefficients with certified accuracy. Computing or estimating the
SVD, prefix masses, Gram matrix, and feasible flows must be charged.
The LP of Theorem~\ref{II:U:thm:flow-lp} instead minimizes a confidence-adjusted
upper bound for the truncated $\ell^1$ closure certificate; it neither
computes nor minimizes \eqref{II:U:eq:cone-defect} directly.

\begin{theorem}[Conditional spectral--conic TV bridge]
\label{II:U:thm:conditional-cone-bridge}
Assume $p=\sum_i\alpha_i f_i$ with the anchored choice $f_1=p$ and $\alpha=e_1$, as well as $\mathcal H_{\rm anchor}(k,N,\mu_{\min},\kappa)$ and
$\mathcal H_{\rm cone}^{(2)}(V,\eta_{\rm cone})$ and choose a
witnessing feasible flow. Decode $\widehat q$ as in
Theorem~\ref{II:U:thm:flows}. It has at most $N$ states, is proper, and
$\mathbb E_{\widehat q}|W|\le V-1$. Define
$t_p(s)=P_p(|W|>s)$ for nonnegative integers $s$, and
$D_T=\sum_{\ell=0}^{T-1}d^\ell$ for $T\ge1$. Then
\begin{align}
 J_\infty
 &\le\sqrt{D_T}\bigl(dV c_{\mathcal U}\epsilon_k
                        +\eta_{\rm cone}\bigr)
       +\frac{2V-1}{\mu_{\min}}t_p(T-1),
       \label{II:U:eq:conditional-cone-j}\\
 \UTV(p,\widehat q)&\le\tfrac12J_\infty.
       \label{II:U:eq:conditional-cone-tv}
\end{align}
The conic defect in this theorem is the $\ell^2$ quantity
\eqref{II:U:eq:cone-defect}, not the total $\ell^1$ defect $J_\infty$.
\end{theorem}
\begin{proof}
For $R_i=x_if_i-z_i\delta_\varepsilon-\sum_{a,j}y_{aij}L_af_j$,
let $A=\sum_i|R_i(\varepsilon)|$. On the branch starting with $a$,
\[
 L_a^*R_i=x_i(L_a^*f_i-F_{\mathcal U}w_{a,i}^{\rm lin})
           +F_{\mathcal U}(x_iw_{a,i}^{\rm lin}-y_{a,i,\cdot}).
\]
Lemma~\ref{II:U:lem:signed-anchor-closure} and
Definition~\ref{II:U:def:cone-defect} give
\[
 \sum_{a,i}\|L_a^*R_i\|_2
 \le dV c_{\mathcal U}\epsilon_k+\eta_{\rm cone}-A.
\]
Cauchy--Schwarz on the $D_T$ suffixes of lengths $0$ through $T-1$
therefore bounds the contribution of $|w|\le T$ by
$\sqrt{D_T}(dV c_{\mathcal U}\epsilon_k+\eta_{\rm cone})$,
since $\sqrt{D_T}\ge1$ absorbs $A$.

The exact residual-tail relation and its useful upper bound are
\begin{equation}
 t_{f_i}(s)=\frac{\sum_{|v|>s}p(u_iv)}{\mu_i}
 \le\frac{t_p(s+|u_i|)}{\mu_i}
 \le\frac{t_p(s)}{\mu_{\min}}.
 \label{II:U:eq:conditional-anchor-tail}
\end{equation}
Positivity bounds the remaining contribution by
\begin{align*}
 \sum_i\sum_{|w|>T}|R_i(w)|
 &\le\sum_i x_it_{f_i}(T)
           +\sum_{a,i,j}y_{aij}t_{f_j}(T-1)\\
 &=\sum_i x_it_{f_i}(T)
           +\sum_j(x_j-\alpha_j)t_{f_j}(T-1)\\
 &\le\frac{Vt_p(T)+(V-1)t_p(T-1)}{\mu_{\min}}\\
 &\le\frac{2V-1}{\mu_{\min}}t_p(T-1).
\end{align*}
Adding head and tail proves \eqref{II:U:eq:conditional-cone-j}.
Since $f_1=p$, Theorem~\ref{II:U:thm:flows} applies with no initial-mixture
error and proves all remaining assertions.
\end{proof}


\paragraph{An alternative direct $\ell^1$ hypothesis.}
One can instead require a feasible flow with
\[
 \eta_1(x,z,y):=\sum_i|x_if_i(\varepsilon)-z_i|
       +\sum_{a,i}\left\|x_iL_a^*f_i-\sum_jy_{aij}f_j\right\|_1
       \le\eta_1^{\rm target}.
\]
The empty word and letter branches have disjoint supports, so
$\eta_1=J_\infty$ exactly and the TV bound is
$\eta_1^{\rm target}/2$, without spectral truncation or a separate
tail term. This directly certifying formulation must not be
identified with Definition~\ref{II:U:def:cone-defect}.
For a KL bound without increasing the number of states, apply the local
smoothing construction in Section~\ref{II:U:sec:flow-smoothing} with a
\emph{prescribed positive upper bound} $\Delta$ on the certified TV.
Its tail assumptions and possible increase in expected visits remain
part of the guarantee. A mean bound alone does not make KL vanish
uniformly, by Theorem~\ref{II:U:thm:kl-fixed-states}.

\subsection{Certified search and the unresolved acquisition theorem}
\label{II:U:sec:acquisition-contract}
An anchor-search heuristic may start with $\varepsilon$, fit a feasible
flow, and add a previously absent child $u_i a$ whose estimated branch
defect is large and whose prefix mass passes a declared threshold.
The exact local identity
\[
 L_a^*f_{u_i}=\frac{P_p(u_i a\Sigma^*)}{P_p(u_i\Sigma^*)}f_{u_i a}
\]
holds whenever the child mass is positive. It closes that branch if its
coefficient is used, but the resulting flows must still jointly satisfy
\eqref{II:U:eq:flow}. Neither the visit budget, the spectral conditioning,
nor all other branch defects are thereby controlled.

\paragraph{What a successful run certifies.}
Freeze the dictionary before the certification sample, or establish a
uniform concentration event covering its adaptive selection. On the event
in Theorem~\ref{II:U:thm:flow-lp}, an actually returned feasible flow with
$\widehat C_L\le2\varepsilon$ certifies
$\UTV(p,\widehat q)\le\varepsilon$. This is a probability-at-least
$1-\delta$ statement when empirical data are used, not an unconditional
deterministic guarantee. The LP may optimize its flows on those data
because its event is uniform over feasible flows for the fixed dictionary.
If successive selected dictionaries are certified on fresh data, allocate
failure probabilities $\delta_r$ with $\sum_r\delta_r\le\delta$.
Reusing a holdout to select dictionaries without such control is not
covered. Numerical flow feasibility and the achieved objective must also
be verified; the unknown optimum alone is not a certificate. This LP
certificate does not require spectral full rank or a condition-number
bound unless the spectral bridge is separately claimed.

\paragraph{The remaining acquisition problem.}
A discovery theorem must specify an independently checkable source class,
an input/access model, and an algorithm, then prove all of the following:
\begin{enumerate}[label=(A\arabic*)]
\item A dictionary of at most $N(\varepsilon)$ anchors is found in
bounded time, with explicitly controlled minimum prefix masses and
spectral conditioning, rather than assumed to exist.
\item A globally feasible flow is obtained with controlled visit budget
$V$ and conic defect. Their dependence on the accuracy must make the
\emph{whole} right side of \eqref{II:U:eq:conditional-cone-j} at most
$2\varepsilon$ for a stated choice of $T$. Merely requiring
$\eta_{\rm cone}\le\varepsilon$ is insufficient when $D_T$, $V$,
$\kappa$, or $\mu_{\min}^{-1}$ grow.
\item Anchor discovery, conditional sampling, spectral/Gram estimation,
finite precision, and optimization have explicit sample, bit, and runtime
costs, including rare-prefix and adaptive-selection effects.
\item The algorithm returns a verified model under the stated assumptions;
outside them, failure to find a certificate is reported as unresolved,
not as nonexistence of a positive approximation.
\end{enumerate}
The smallness of $\sigma_{k+1}(H_p)$ alone proves none of these
acquisition claims. Proposition~\ref{II:U:prop:positive-obstruction} rules out
an exact universal positive-realization implication even at zero spectral
tail. It does not rule out approximation with a growing state count:
conditioning $p$ on $|W|\le T$ gives a finite-support law at TV distance
$t_p(T)$, realizable by a finite prefix tree. This construction neither
provides a small state bound nor finite KL for unbounded-support $p$.
The transfer-operator result in Subsection~\ref{III:U:sec:oracle} concerns
already supplied contracting two-state homographic models and supplies
no missing anchor-discovery theorem.


\subsection{What can be certified over a finite candidate family}
\label{II:U:sec:finite-family-search}
Existence of a good hidden anchor tree does not prove that a local greedy
rule reaches it within the same number of nodes. A finite exhaustive
search has the following weaker, explicit guarantee.

\begin{theorem}[Finite-family search with a common validation event]
\label{II:U:thm:finite-family-search}
Fix $M$ candidate prefix laws $f_1=p,f_2,\ldots,f_M$ before validation,
and a finite family $\mathcal D$ of subsets containing index one, each
of size at most $N_*$. Assume conditional samples from these laws are
available. Each candidate dictionary is started with $\alpha=e_1$, so $p=\sum_i\alpha_i f_i$ holds exactly. For an integer horizon $L\ge1$, take $n_i$ independent samples per candidate
law, with no independence required between different candidate laws,
and use
\[
 b_i=\min\left\{2,\sqrt{
    \frac{2(K_L\log2+\log(M/\delta))}{n_i}}\right\},\qquad
 K_L=1+\sum_{j=0}^Ld^j.
\]
Solve the flow LP for every dictionary in $\mathcal D$ with budget $V$,
using the same histograms and the applicable $b_i$. With probability
at least $1-\delta$, every feasible returned model satisfies its
certificate $\UTV(p,q)\le\widehat C_L/2$ simultaneously.
If one candidate dictionary admits a flow satisfying
\begin{equation}
 J_\infty+2\sum_i x_it_i(L)+2S(x)\le2\varepsilon,
 \label{II:U:eq:finite-family-witness}
\end{equation}
then, with exact LP solves, the search finds a certified model of error
at most $\varepsilon$ after at most $|\mathcal D|$ solves.
Otherwise it may return \texttt{UNRESOLVED}. For all subsets of size
at most $N_*$, the number of solves is bounded by
$\sum_{j=0}^{N_*-1}\binom{M-1}{j}$, with terms beyond $M-1$ omitted.
\end{theorem}
\begin{proof}
Apply the histogram concentration argument of
Theorem~\ref{II:U:thm:flow-lp} once to the $M$ fixed candidate laws.
Its simultaneous event covers every feasible flow on every subset, so
subsequent dictionary selection does not require another union bound.
For the witnessing flow,
$\widehat C_L^*\le J_L+2S(x)
\le J_\infty+2\sum_i x_it_i(L)+2S(x)$.
Condition \eqref{II:U:eq:finite-family-witness} therefore forces an accepted
dictionary when it is examined. The finite subset count is immediate.
\end{proof}


This theorem neither discovers the candidate universe nor promises a
small search. Access to conditional samples, certification of prefix
masses, numerical LP feasibility, and spectral data are separate costs.
For numerical LP solves, a verified objective gap must be included in
\eqref{II:U:eq:finite-family-witness}. An existential witness-tree assumption
does not yield the proposed greedy guarantee. In particular, ordinary
column pivoting is not the strong rank-revealing QR algorithm of
Gu--Eisenstat \citep{gu1996}; moreover the inequality
$2/\sqrt{k(N_*-k)+1}\ge1$ is already false at $(k,N_*)=(3,6)$.

\subsection{Positive closure in the correct dual state coordinates}
\label{II:U:sec:dual-positive-closure}
For $p(w)=c^\top A_wb$, prefix residuals are represented by
\emph{dual} state vectors. The vectors $A_ub$ instead represent suffix
columns and cannot be substituted without changing the word order.

\begin{theorem}[A sufficient dual-coordinate positive closure condition]
\label{II:U:thm:dual-positive-closure}
Let $p$ have a real minimal mean-square-stable realization of dimension
$k$. Choose anchors with $u_1=\varepsilon$, masses
$\mu_i\ge\mu_{\min}>0$, and vectors
\[
 \ell_i=\mu_i^{-1}A_{u_i}^\top c,\qquad
 (\mathcal R_b\ell)(v)=\ell^\top A_vb.
\]
Assume the $\ell_i$ span $\mathbb R^k$. Suppose nonnegative numbers
$t_{aij}$ satisfy, for every $i,a$,
\begin{align}
 \sum_{a,j}t_{aij}&=1-f_i(\varepsilon),\qquad
 f_i(\varepsilon)\ge\gamma_{\rm stop}>0,\notag\\
 \left\|A_a^\top\ell_i-\sum_jt_{aij}\ell_j\right\|_2
 &\le\eta_{\rm pos}.\label{II:U:eq:dual-positive-hypothesis}
\end{align}
Then the conditioned-anchor hypothesis holds with
$\kappa=\|(V_k^*F_{\mathcal U})^\dagger\|<\infty$.
Moreover, there is a feasible flow with
\[
 V=\gamma_{\rm stop}^{-1},\qquad
 \eta_2\le dV\|\mathcal R_b\|\eta_{\rm pos},\qquad
 \|\mathcal R_b\|^2=\|P_p\|.
\]
Thus Theorem~\ref{II:U:thm:conditional-cone-bridge} applies with
$\epsilon_k=0$. This result presupposes the supplied spanning anchors
and the jointly normalized positive coefficients; it is not an
anchor-discovery theorem.
\end{theorem}
\begin{proof}
The correct identities are
$f_i=\mathcal R_b\ell_i$ and
$L_a^*f_i=\mathcal R_bA_a^\top\ell_i$.
Minimality makes $P_p\succ0$, so $\mathcal R_b$ is injective and
the spanning assumption makes $F_{\mathcal U}$ have rank $k$.
Its image is the right singular space of $H_p$; hence the stated
$\kappa$ is finite and the signed closure in
Lemma~\ref{II:U:lem:signed-anchor-closure} is exact.

Take $\alpha=e_1^\top$, put $T_a=(t_{aij})$ and $T=\sum_aT_a$. Its row sums are at most
$1-\gamma_{\rm stop}$. Therefore
$x=\alpha(I-T)^{-1}\ge0$ and $\sum_i x_i\le V$.
Set $y_{aij}=x_it_{aij}$ and $z_i=x_if_i(\varepsilon)$.
These satisfy both flow equations. Their stop defect is zero, and
the exact signed closure gives
\[
 \eta_2=\sum_{a,i}x_i
 \left\|L_a^*f_i-\sum_jt_{aij}f_j\right\|_2
 \le dV\|\mathcal R_b\|\eta_{\rm pos}.
\]
Finally $\|\mathcal R_b\ell\|_2^2=\ell^\top P_p\ell$ proves
the asserted operator norm. No assertion that this cone is simplicial
is needed; it may have more than $k$ generators.
\end{proof}


\subsection{Conditional predictions and the effect of smoothing}
\label{II:U:sec:conditional-event-bounds}
Let $p_u$ denote the next-event law on $\Sigma\cup\{\mathrm{STOP}\}$
after observing a prefix $u$. All KL quantities below use natural logs.

\begin{theorem}[Prefix-conditioned TV and a valid smoothing bound]
\label{II:U:thm:conditional-event-bounds}
Let $p,q$ be proper word laws with $\UTV(p,q)\le\varepsilon$.
For $s=P_p(u\Sigma^*)\ge\mu_{\rm pref}>\varepsilon$,
$P_q(u\Sigma^*)>0$ and
\begin{equation}
 \UTV(p_u,q_u)\le\frac{\varepsilon}{s}
                         \le\frac{\varepsilon}{\mu_{\rm pref}}.
 \label{II:U:eq:conditional-event-tv}
\end{equation}
Suppose $q$ is an $N$-state automaton with
$1+\mathbb E_q|W|\le V$, and apply the state-preserving local smoothing
of Section~\ref{II:U:sec:flow-smoothing} with parameter $0<\Delta<1$.
Put $e=\varepsilon+\Delta V$. If $\mu_{\rm pref}>e$, then
\begin{align}
 \UTV(p_u,(q_\Delta)_u)&\le e/\mu_{\rm pref},\notag\\
 D_{\rm KL}(p_u\|(q_\Delta)_u)
 &\le\min\left\{\log\frac{d+1}{\Delta},\;
 \frac{2(d+1)}{\Delta}\left(\frac e{\mu_{\rm pref}}\right)^2
 \right\}.\label{II:U:eq:conditional-event-kl}
\end{align}
The number of states remains $N$; no preservation of the visit budget
after smoothing is asserted.
\end{theorem}
\begin{proof}
Let $E=u\Sigma^*$, $t=Q(E)$, and
$A=\sum_{w\in E}|p(w)-q(w)|$. The complement has total variation
mass at least $|s-t|$, so $A+|s-t|\le2\varepsilon$.
Also $t\ge s-\varepsilon>0$. Consequently
\[
 \sum_{w\in E}\left|\frac{p(w)}s-\frac{q(w)}t\right|
 \le\frac{A+|s-t|}{s}\le\frac{2\varepsilon}s.
\]
Mapping a conditional suffix to its first letter or STOP cannot
increase TV; this proves \eqref{II:U:eq:conditional-event-tv}.
The sum of absolute differences of next-event masses is \emph{at most},
not necessarily equal to, the sum over all completed suffixes.

Couple the original and smoothed automata until a smoothing choice
occurs. For an original word $W$, there are $|W|+1$ event choices, so
$\UTV(q,q_\Delta)\le
\mathbb E[1-(1-\Delta)^{|W|+1}]\le\Delta V$.
Apply the first part to $p,q_\Delta$. Every next-event probability
under $q_\Delta$, for any posterior state mixture, is at least
$a=\Delta/(d+1)$. For probability vectors $r,z$ with $z_j\ge a$,
the inequality $\log t\le t-1$ gives
\[
 D_{\rm KL}(r\|z)\le\sum_j\frac{(r_j-z_j)^2}{z_j}
 \le\frac{2\UTV(r,z)^2}{a}.
\]
The last step uses that positive and negative differences each sum to
$\UTV(r,z)$. Also $D_{\rm KL}(r\|z)\le\log(1/a)$ by dropping
the nonnegative entropy of $r$. These prove the two KL bounds.
\end{proof}


\paragraph{Counterexample to prefix-independent smoothing error.}
Take $p=q$ on two states with initial row $(1/2,1/2)$,
\[
 T_a=\Udiag(t,0),\quad T_b=\Udiag(0,1/2),\quad
 \tau=(1-t,1/2)^\top,
 \qquad t=10^{-6},\quad\Delta=10^{-3}.
\]
After prefix $a$, the unsmoothed next-event law (STOP,$a$,$b$) is
$(1-t,t,0)$, and its prefix mass is $t/2>0$. After local smoothing,
the posterior weight of the first state is
\[
 \omega=\frac{(1-\Delta)t+\Delta/3}
                   {(1-\Delta)t+2\Delta/3}.
\]
The new next-event law is approximately
$(0.7499565235,0.0003338336,0.2497096429)$.
Its KL from the original law exceeds $0.2877$, whereas a proposed uniform smoothing bound at $\varepsilon=0$ is
$\Delta\log(3/\Delta)<0.0081$. The changed posterior after a rare
prefix cannot be ignored. This example also has original visit budget
less than two; it is not a failure caused by infinite expected length.


\subsection{Finite positive states need not be finite prefix anchors}
\label{II:U:sec:rank-two-obstruction}
The rank-two results below are conditional on a supplied, exact
two-dimensional model and a certified interval of probability laws.
They do not complete arbitrary-rank spectral learning. They also
distinguish latent generators, limiting residuals, and actual finite
prefix residuals; these are different admissible dictionaries.

\begin{lemma}[The aggregate of a minimal proper-word realization is stable]
\label{II:U:lem:proper-aggregate-stable}
Let $p(w)=c^\top A_wb$ be a minimal finite-dimensional real realization
of a probability law on finite words, and put $T=\sum_a A_a$.
Then $T^n\to0$ and $\rho(T)<1$.
\end{lemma}
\begin{proof}
Choose words $u_i,v_j$ so that the matrix $O$ with rows
$c^\top A_{u_i}$ and the matrix $C$ with columns $A_{v_j}b$ are
square and invertible; minimality guarantees these choices. For each
fixed pair $(i,j)$,
\[
 (OT^nC)_{ij}=\sum_{|w|=n}p(u_iwv_j)
 \le P_p\{|W|=n+|u_i|+|v_j|\}\longrightarrow0.
\]
Hence $T^n=O^{-1}(OT^nC)C^{-1}\to0$. In finite dimension this
implies that every eigenvalue has modulus strictly below one.
\end{proof}

Thus a Perron eigenvalue equal to one for the unnormalized aggregate
is incompatible with this convention. A normalized continuation
operator is a different object.

\begin{proposition}[A strictly positive two-state prefix obstruction]
\label{II:U:prop:golden-prefix-obstruction}
Set
\[
 A_0=\frac1{10}\begin{pmatrix}2&1\\1&1\end{pmatrix},\quad
 A_1=\frac1{10}\begin{pmatrix}1&1\\1&2\end{pmatrix},\quad
 \alpha=(1/2,1/2),\quad \tau=(1/2,1/2)^\top.
\]
The law $p(w)=\alpha A_w\tau$ is proper and has Hankel rank two.
Both letters map the positive orthant strictly into its interior,
but no finite dictionary of actual prefix residuals containing $p$
in its convex hull has exact positive closure.
\end{proposition}
\begin{proof}
Here $T\mathbf1=\frac12\mathbf1$, so termination is geometric.
The two-row reachability matrix with rows $\alpha,\alpha A_0$
and the two-column observability matrix with columns $\tau,A_0\tau$
both have determinant $-1/40$, proving rank two.
For this example write $f_t(v)=(t,1-t)A_v\tau$; its first coordinate
updates by
\[
 F_0(t)=\frac{1+t}{2+t},\qquad F_1(t)=\frac1{3-t}.
\]
On $[0,1]$ their images are $[1/2,2/3]$ and $[1/3,1/2]$.
With $r=(\sqrt5-1)/2$, both maps preserve $(1-r,r)$, and every
finite-prefix coordinate lies in that interval. The words $0^n$
and $1^n$ approach its two endpoints. For any finite dictionary let
$m<r$ be its maximum coordinate. Then
\[
 F_0(m)-m=\frac{1-m-m^2}{2+m}>0.
\]
The next residual is outside the dictionary's convex hull. The
coordinate map is affine and injective since
$f_t(0)=1/10+t/20$. Nonnegative representations of normalized laws
have coefficient sum one, so convex and conic closure give the same
obstruction after normalizing each nonzero shift.
\end{proof}


The same obstruction rules out $J_\infty=0$ for every finite feasible
occupation flow with this source in its initial mixture: restrict to
states with positive occupation. Zero-occupation states receive no
flow; zero defect would make the active dictionary invariant. For a
fixed finite visit budget the feasible flow set is compact, hence a
nonempty such problem has strictly positive minimum defect. No
uniform positive lower bound over increasing dictionaries is claimed.
The distinction between positive rational realization and finite
generation by residuals is classical \citep{denis2008}.

For example,
\[
 t_\varepsilon=\tfrac12,\quad t_0=\tfrac35,\quad t_{00}=\tfrac8{13},
 \qquad f_{00}=-\tfrac2{13}f_\varepsilon+\tfrac{15}{13}f_0.
\]
Thus linear QR residuals vanish after two independent anchors while
conic residuals need not. Also $2T$ has eigenvalues $1,1/5$, yet
$0^n$ and $1^n$ have different limiting conditional states. A gap for
the normalized aggregate does not give convergence of every letter
product to a common ray.

\begin{lemma}[A sufficient simplex-expanding exchange]
\label{II:U:lem:simplex-exchange}
Let $f_1,\ldots,f_R$ be linearly independent normalized laws, and let
another law satisfy $f=\sum_i\beta_i f_i$, where $\sum_i\beta_i=1$.
If $\beta_j>1$ for one index $j$ and $\beta_i\le0$ for every $i\ne j$,
replacing $f_j$ by $f$ preserves linear independence and strictly
enlarges the convex hull.
\end{lemma}
\begin{proof}
The identity
\[
 f_j=\beta_j^{-1}f+
       \sum_{i\ne j}(-\beta_i/\beta_j)f_i
\]
is a convex combination. Independence follows from $\beta_j\ne0$,
and $f$ was outside the old simplex by uniqueness of its coordinates.
\end{proof}

This is a sufficient exchange test, not a termination theorem.
In higher affine dimension an exterior point can have several
positive coordinates, so an eligible exchange need not exist.

\subsection{Rank-two monotone realization and finite-prefix approximation}
\label{II:U:sec:rank-two-monotone}
Choose two distinct normalized residual laws $f_0,f_1$ spanning the
rank-two residual space. In this subsection the convention is
\[
 f_t=(1-t)f_0+tf_1,\qquad
 L_1=\|f_1-f_0\|_1,\qquad \|f_t-f_s\|_1=L_1|t-s|.
\]
The notation differs from the first-coordinate convention used in
Proposition~\ref{II:U:prop:golden-prefix-obstruction}; changing orientation
does not change any law or certificate. There is a real realization
$f_t(v)=(1-t,t)A_v\tau$. By Lemma~\ref{II:U:lem:proper-aggregate-stable}
and normalization, $(I-\sum_a A_a)^{-1}\tau=\mathbf1$.
For $A_a=\left(\begin{smallmatrix}a_a&b_a\\c_a&d_a\end{smallmatrix}\right)$
define
\begin{align}
 w_a(t)&=(1-t,t)A_a\mathbf1,\qquad s(t)=f_t(\varepsilon)
                  =1-\sum_a w_a(t),\notag\\
 F_a(t)&=\frac{b_a+(d_a-b_a)t}{w_a(t)},\qquad
 F_a'(t)=\frac{\det A_a}{w_a(t)^2}.
 \label{II:U:eq:rank-two-projective}
\end{align}
When $w_a(t)>0$, $L_a^*f_t=w_a(t)f_{F_a(t)}$.

\begin{theorem}[Monotone rank-two limiting and prefix realizations]
\label{II:U:thm:rank-two-monotone}
Suppose a compact interval $I$ contains $t_0$ with $p=f_{t_0}$,
every $f_t$ for $t\in I$ is a probability law, and every letter satisfies
$w_a(t)>0$, $F_a(I)\subseteq I$, and
$0<F_a'(t)\le\bar\gamma_a<1$ on $I$.
Let $r_a$ be its unique fixed point and set
\[
 r_-=\min_a r_a=r_{a_-},\qquad
 r_+=\max_a r_a=r_{a_+}.
\]
Assume $r_-<t_0<r_+$ and
$s_*=\min_{[r_-,r_+]}s(t)>0$.
Then the closed convex hull of all prefix residuals is
$\operatorname{conv}(f_{r_-},f_{r_+})$. In this ordered dictionary,
the following matrices give an exact proper two-state realization:
\begin{align}
 \lambda_a^\pm&=\frac{F_a(r_\pm)-r_-}{r_+-r_-},\qquad
 B_a=\begin{pmatrix}
 w_a(r_-)(1-\lambda_a^-)&w_a(r_-)\lambda_a^-\\
 w_a(r_+)(1-\lambda_a^+)&w_a(r_+)\lambda_a^+
 \end{pmatrix},\notag\\
 \alpha&=\left(\frac{r_+-t_0}{r_+-r_-},
              \frac{t_0-r_-}{r_+-r_-}\right),\qquad
 \tau_B=(s(r_-),s(r_+))^\top.
 \label{II:U:eq:rank-two-limit-model}
\end{align}

For $L\ge1$, let
$l_L=F_{a_-}^L(t_0)$ and $h_L=F_{a_+}^L(t_0)$.
They are actual prefix coordinates, with
$r_-<l_L<t_0<h_L<r_+$.
Replace $r_-,r_+$ in the initial and stopping vectors by $l_L,h_L$,
and replace each successor $F_a(t)$ in the transition formula by
$\Pi_{[l_L,h_L]}F_a(t)$, for $t=l_L,h_L$. Denote the resulting
matrices by $\widehat B_{a,L}$, the law by $q_L$, and
$Q_L=\sum_a\widehat B_{a,L}$. Then
\[
 x_L=\alpha_L(I-Q_L)^{-1}\ge0,\qquad x_L\mathbf1\le s_*^{-1},
\]
and its feasible occupation flows have exact defect
\begin{equation}
 J_L^{\rm clip}
 =L_1\sum_{t\in\{l_L,h_L\}}x_{L,t}
        \sum_a w_a(t)\operatorname{dist}(F_a(t),[l_L,h_L]).
 \label{II:U:eq:rank-two-clip-defect}
\end{equation}
Here $J_L^{\rm clip}$ is an untruncated $J_\infty$ certificate;
it is not the finite-word objective $J_L$ of Section~\ref{II:U:sec:flow-lp}.
Writing $\delta_L=\max(l_L-r_-,r_+-h_L)$ gives
\begin{equation}
 \UTV(p,q_L)\le\tfrac12J_L^{\rm clip}
 \le\frac{L_1(1-s_*)}{2s_*}\delta_L.
 \label{II:U:eq:rank-two-clip-tv}
\end{equation}
No pointwise ordering of the different maps $F_a$ is required.
\end{theorem}
\begin{proof}
For an increasing contraction with fixed point $r_a$, a point below
$r_a$ moves strictly toward $r_a$ without crossing it, and a point
above does the same from the other side. Thus each map preserves
$[r_-,r_+]$. All reachable points lie in it; the repeated letters
$a_-^n,a_+^n$ approach its endpoints. Affinity and continuity into
$\ell^1$ prove the closed convex hull assertion.

All $\lambda_a^\pm$ lie in $[0,1]$. The rows sum to $w_a(r_\pm)$,
and the residual recurrences hold exactly. Alternatively, with
$S=\left(\begin{smallmatrix}1-r_-&r_-\\1-r_+&r_+\end{smallmatrix}\right)$,
$B_a=SA_aS^{-1}$ and $\tau_B=S\tau$, which proves equality of
every word coefficient. Each stop probability is at least $s_*$,
so the model is proper.

The same row normalization holds after clipping, with stop at least
$s_*$. Hence $\rho(Q_L)\le1-s_*$ and the displayed occupation vector
exists. Take $y_{aij}=x_i\widehat B_{a,L}(i,j)$ and
$z_i=x_i\tau_{L,i}$ to obtain the flow equations. The source is
exactly the mixture $\alpha_L(f_{l_L},f_{h_L})^\top$.
The stop residual is zero, and the residuals for different first
letters have disjoint support. Each has norm
$L_1w_a(t)|F_a(t)-\Pi_{[l_L,h_L]}F_a(t)|$, proving
\eqref{II:U:eq:rank-two-clip-defect}.
Invariance bounds every distance by $\delta_L$; also
$\sum_a w_a(t)\le1-s_*$ and $x_L\mathbf1\le1/s_*$.
Apply Theorem~\ref{II:U:thm:flows}.
\end{proof}


\begin{corollary}[Exact cross-ratio depth rates]
\label{II:U:cor:rank-two-cross-ratio}
For either extremal map let $r$ be its attracting fixed point and
$\gamma=F'(r)\in(0,1)$. If its second fixed point $\bar r$ is finite,
put $c=(t_0-r)/(t_0-\bar r)$. Then
\begin{equation}
 F^L(t_0)=\frac{r-c\gamma^L\bar r}{1-c\gamma^L},
 \qquad
 |F^L(t_0)-r|
 =\frac{|r-\bar r|\,|c|\,\gamma^L}{|1-c\gamma^L|}.
 \label{II:U:eq:rank-two-cross-ratio}
\end{equation}
If $F$ is affine, the formula is instead
$F^L(t_0)-r=\gamma^L(t_0-r)$.
Consequently $\delta_L\le C\gamma_*^L$, with
$\gamma_*=\max(F_{a_-}'(r_-),F_{a_+}'(r_+))<1$
and an explicit finite, source-dependent constant $C$.
\end{corollary}
\begin{proof}
The two fixed points are distinct because $F'(r)<1$.
For a nonaffine map the second fixed point is real and outside $I$,
since the contraction has only one fixed point in $I$.
The map $\psi(t)=(t-r)/(t-\bar r)$ conjugates $F$ to multiplication
by $F'(r)$, giving \eqref{II:U:eq:rank-two-cross-ratio}.
Since $t_0,r$ lie on the same side of $\bar r$, $c<1$ and
$1-c\gamma^L\ge\min(1,1-c)>0$. One may use
$C_r=|r-\bar r||c|/\min(1,1-c)$, or $C_r=|t_0-r|$
in the affine case, and take the larger endpoint constant.
\end{proof}

The closed interval here is a convex hull, not necessarily the IFS
attractor. A local multiplier $\gamma_*$ is justified by the exact
cross-ratio identity, not by treating it as a global Lipschitz constant.

\subsection{Both orientations: endpoint acquisition without word enumeration}
\label{II:U:sec:rank-two-envelope}
The even/odd powers of one decreasing letter do not generally find the
extrema created by mixed words. For instance, the proper rank-two model
\[
 A_+=\tfrac1{10}\begin{pmatrix}2&1\\1&1\end{pmatrix},\quad
 A_-=\tfrac1{20}\begin{pmatrix}1&2\\2&1\end{pmatrix},\quad
 \alpha=(1/2,1/2),\quad \tau=(11/20,13/20)^\top
\]
has first-coordinate updates
$F_+(t)=(1+t)/(2+t)$ and $F_-(t)=(2-t)/3$.
All powers of $F_-$ fix $1/2$, but
$F_-(F_+(1/2))=7/15$. The following procedure treats all letters.

\begin{theorem}[Rank-two endpoint iteration for either orientation]
\label{II:U:thm:rank-two-envelope}
Use the affine-law setup above. Suppose every $f_t$ for $t$ in a
compact interval $I$ is a probability law, $t_0\in I$,
$w_a(t)>0$, $F_a(I)\subseteq I$, and every $F_a$ is monotone
(increasing, decreasing, or constant) and $\gamma$-Lipschitz on $I$,
with $0<\gamma<1$. Assume $s(t)\ge s_*>0$ on $I$.
Initialize $K_0=\{t_0\}$ and iterate
\begin{equation}
 K_{n+1}=\operatorname{conv}
   \left(\{t_0\}\cup\bigcup_a F_a(K_n)\right)
       =[l_{n+1},h_{n+1}].
 \label{II:U:eq:rank-two-envelope}
\end{equation}
Then $K_n$ is the convex hull of coordinates of prefixes of length
at most $n$, and its endpoints have actual prefix witnesses.
The intervals increase to a unique compact fixed interval $K_*$
whose endpoints are limits of prefix residuals, with
\[
 d_H(K_n,K_*)\le\gamma^n\operatorname{diam}I.
\]
If the reachable residual space has rank two, $K_*$ is nondegenerate,
and $K_n$ is nondegenerate for all sufficiently large $n$.
The endpoint laws of $K_*$ give an exact two-state realization.
Clipping into any nondegenerate $K_n$ gives a proper two-state law with
\[
 \UTV(p,q_n)\le
 \frac{L_1(1-s_*)}{2s_*}\gamma^n\operatorname{diam}I.
\]
At each iteration only the $2d$ endpoint images and $t_0$ are needed.
Thus $n$ iterations require $O(dn)$ exact real-arithmetic evaluations
and comparisons; no bit-complexity bound is asserted.
\end{theorem}
\begin{proof}
A continuous monotone map sends an interval to the interval between
its endpoint images. Induction therefore identifies $K_n$ with the
convex hull of the finite set of coordinates at depths at most $n$;
each new extreme is attained by a previous extreme followed by one
letter, or by the empty prefix. The interval map in
\eqref{II:U:eq:rank-two-envelope} is monotone for inclusion and is a
$\gamma$-contraction in Hausdorff distance: each image map is, and
adjoining the same point and taking interval hulls cannot increase the
maximum endpoint discrepancy. The space of nonempty compact
subintervals of $I$ is complete. The contraction theorem gives a
unique $K_*$ and the error bound. Since $K_0\subseteq K_1$, induction
gives nesting and identifies $K_*$ as the closed convex hull of all
reachable coordinates.

Rank two rules out a singleton reachable hull. Its endpoint laws are
probabilities by continuity and are invariant under all letters.
The barycentric construction and clipping proof of
Theorem~\ref{II:U:thm:rank-two-monotone} now apply on $K_*$ and $K_n$.
Prefix witnesses can be retained as endpoint backpointers; there are
at most two new witnesses per round. Exact arithmetic and certified
input hypotheses, including the interval of nonnegative laws, are
assumptions of this construction.
\end{proof}

This is a constructive theorem for a declared rank-two interval class.
It does not discover that class from noisy Hankel data, certify its
positivity interval, or establish the missing arbitrary-rank acquisition
theorem. General contraction alone gives no fixed-$R$ simplex theorem
in higher dimension.

\subsection{An exact golden-ratio clipping certificate}
\label{II:U:sec:golden-constructive}
Return to the first-coordinate convention of
Proposition~\ref{II:U:prop:golden-prefix-obstruction}. In the ordered
limiting dictionary $(f_r,f_{1-r})$ the exact matrices are
\[
 B_0=\begin{pmatrix}(3+\sqrt5)/20&0\\1/5&(3-\sqrt5)/20\end{pmatrix},
 \qquad
 B_1=\begin{pmatrix}(3-\sqrt5)/20&1/5\\0&(3+\sqrt5)/20\end{pmatrix}.
\]
The initial vector remains $(1/2,1/2)$ and both stops are $1/2$.
The matrices equal $SA_aS^{-1}$ for
$S=\left(\begin{smallmatrix}r&1-r\\1-r&r\end{smallmatrix}\right)$.
The limiting laws are not finite-prefix residuals.

\begin{proposition}[Exact two-prefix defect on the golden-ratio source]
\label{II:U:prop:golden-clip}
For $L\ge1$ set
\[
 a=F_0^L(1/2),\quad
 \gamma=\frac{7-3\sqrt5}{2},\quad h=9-4\sqrt5,\quad z_L=h\gamma^L.
\]
In the ordered dictionary $(f_a,f_{1-a})$, clipping gives
\[
 \widehat B_{0,L}=
 \begin{pmatrix}
 (2+a)/10&0\\
 (-a^2+3a-1)/(10(2a-1))&(-a^2+4a-2)/(10(2a-1))
 \end{pmatrix}.
\]
\[
 \widehat B_{1,L}=P\widehat B_{0,L}P,\qquad
 P=\begin{pmatrix}0&1\\1&0\end{pmatrix}.
\]
These matrices are nonnegative, have total transition row sum $1/2$,
and yield occupation vector $x_L=(1,1)$ from $(1/2,1/2)$.
Their exact untruncated flow defect is
\begin{equation}
 J_L^{\rm clip}
 =\frac{1-a-a^2}{25}
 =\frac{z_L}{5(1+z_L)^2},\qquad
 \UTV(p,q_L)\le\frac{z_L}{10(1+z_L)^2}.
 \label{II:U:eq:golden-exact-defect}
\end{equation}
\end{proposition}
\begin{proof}
Here $3/5\le a<r$. The inward successors lie in $[1-a,a]$;
the two outward ones are clipped to their respective endpoints.
Their barycentric coefficients give the displayed matrices.
The aggregate has both row and column sums $1/2$, hence $x_L=(1,1)$.
If $f^{[1]},f^{[2]}$ are the laws from the two original pure states,
then $(1,-1)A_0=(1/10,0)$ and $(1,-1)A_1=(0,-1/10)$.
Their difference has positive mass $1/10$ on words starting with $0$
and negative mass $-1/10$ on words starting with $1$, with zero
empty-word difference. Thus
$\|f_t-f_s\|_1=|t-s|/5$.
The outward coordinate displacement is
$(1-a-a^2)/(2+a)$ and its letter probability is $(2+a)/10$.
The two occupied rows therefore give $(1-a-a^2)/25$.
The cross-ratio formula gives
$r-a=\sqrt5\,z_L/(1+z_L)$, so
$1-a-a^2=5z_L/(1+z_L)^2$.
Theorem~\ref{II:U:thm:flows} completes the proof.
\end{proof}

At $L=1$, $J_1^{\rm clip}=1/625$ and $\UTV(p,q_1)\le1/1250$.
These are certificate values, not assertions that the actual TV
distance or the optimized flow objective equals the displayed bound.


\subsection{Observable coordinates, explicit rounding, and conditional sampling}
\label{II:U:sec:rank-two-numerical}
Fix a nondegenerate two-prefix dictionary $(f_l,f_h)$ from either
rank-two construction, with $l<h$ and $l\le t_0\le h$.
The exact clipped model is denoted by $q$, with stop probabilities at
least $s_*>0$ and defect $J^{\rm clip}$. Choose a fixed measurable
suffix event $E\subseteq\Sigma^*$ such that
$c_h=f_h(E)>c_l=f_l(E)$, replacing $E$ by its complement if necessary.
Put $\Delta_E=c_h-c_l>0$. For any coordinate $t$,
\[
 \frac{t-l}{h-l}=\frac{f_t(E)-c_l}{c_h-c_l}.
\]
If $L_E=|f_{r_+}(E)-f_{r_-}(E)|$ denotes separation of the limiting
laws, then the normalization is
\begin{equation}
 \Delta_E=L_E\,\frac{h-l}{r_+-r_-},\qquad
 D_{l,h}:=\|f_h-f_l\|_1=L_1(h-l).
 \label{II:U:eq:observable-gap-correct}
\end{equation}
The interval width cannot be silently omitted. The transition
barycentric weights are the clipped ratios for the successor
residuals $f^{(u_l a)},f^{(u_h a)}$, and the initial weight is the
ratio for $p$. Thus these quantities are observable from the fixed
query set
\[
 \mathcal W=\{\varepsilon,u_l,u_h\}
       \cup\{u_l a,u_h a:a\in\Sigma\},\qquad M=|\mathcal W|\le3+2d.
\]
All these prefixes have positive mass under the interval hypotheses.

\begin{lemma}[Stopped-automaton perturbation by coupling]
\label{II:U:lem:stopped-coupling}
Let $q,\widetilde q$ be proper finite-state word laws with initial
vectors $\alpha,\widetilde\alpha$ on the same states. Let $K_i$ and
$\widetilde K_i$ be the state-$i$ probability vectors on the finite
outcomes $\{\mathrm{STOP}\}\cup(\Sigma\times\{1,\ldots,N\})$.
If the visit vector $\nu=\sum_{n\ge0}\alpha(\sum_a B_a)^n$ is finite, then
\[
 \UTV(q,\widetilde q)
 \le\UTV(\alpha,\widetilde\alpha)
       +\sum_i\nu_i\UTV(K_i,\widetilde K_i).
\]
\end{lemma}
\begin{proof}
Maximally couple initial states and, while the states agree, maximally
couple the next finite outcome. The words agree unless a mismatch
occurs. At any visit to $i$ before mismatch, its probability is at
most $\UTV(K_i,\widetilde K_i)$. Summing these probabilities is bounded
by the expected number $\nu_i$ of visits in the first automaton.
The coupling inequality gives the assertion. This proof uses the
original visits, not visits of the perturbed chain.
\end{proof}


\begin{theorem}[A deterministic rank-two estimation and rounding certificate]
\label{II:U:thm:rank-two-rounding}
Suppose every queried event probability $f^{(u)}(E)$, $u\in\mathcal W$,
has an estimate with absolute error at most $\epsilon_E<\Delta_E/4$.
At each of the two anchors suppose the estimated vector of stop and
letter probabilities is nonnegative, sums to one, and has $\ell^1$
error at most $\epsilon_m<s_*$. Define
\[
 e=\frac{4\epsilon_E}{\Delta_E},\qquad
 h_b=2^{-b},\quad b\in\mathbb N,\ b\ge1.
\]
Form initial and successor barycentric weights from the estimated
ratios, clipped to $[0,1]$. Multiply each estimated letter mass by
its two barycentric weights to obtain a preliminary transition row.
Round each of its $2d$ transition entries down to a multiple of $h_b$,
and assign the remaining row mass to STOP. Round the initial weight
of the second state down to the same grid and define the first as
its complement. The resulting two-state law $\widetilde q_b$ is proper
and satisfies
\begin{equation}
 \UTV(p,\widetilde q_b)\le
 \frac{J^{\rm clip}}2+e+h_b+
 \frac{\epsilon_m/2+e+2dh_b}{s_*}.
 \label{II:U:eq:rank-two-rounding}
\end{equation}
No perturbation of the original occupation weights is assumed.
\end{theorem}
\begin{proof}
Let $a=c_l,b'=c_h$ and let $z$ denote any queried event probability.
If $z\in[a,b']$, write $z=(1-t)a+tb'$ with $t\in[0,1]$.
The numerator of the ratio perturbation, after multiplication by the
estimated denominator, is
$\delta z-(1-t)\delta a-t\delta b'$, of absolute value at most
$2\epsilon_E$. The denominator is at least
$\Delta_E-2\epsilon_E$. Clipping cannot increase the discrepancy.
If $z<a$ or $z>b'$, the same bound follows by comparison with the
corresponding endpoint. Thus every clipped barycentric error,
including the initial error, is at most
$2\epsilon_E/(\Delta_E-2\epsilon_E)\le e$.

For a preliminary transition row, separate the error in stop/letter
masses from the error in splitting each letter between two states.
The former has TV at most $\epsilon_m/2$ and the latter at most
$\sum_a\widehat w_a e\le e$. Downward rounding transfers less than
$2dh_b$ mass from transitions to STOP, hence changes row TV by at most
$2dh_b$. The initial TV error is at most $e+h_b$.
Because the stop/letter vectors are normalized, their stop-coordinate
error is at most $\epsilon_m/2$. Downward rounding only increases the
stop mass, so every final stop is at least
$s_*-\epsilon_m/2>0$. Properness follows.
Apply Lemma~\ref{II:U:lem:stopped-coupling}, using the original visit bound
$\nu\mathbf1\le1/s_*$, and then
$\UTV(p,q)\le J^{\rm clip}/2$.
\end{proof}

The theorem concerns represented dyadic transition probabilities and
their mathematical law. Runtime sampling bias and unaccounted
floating-point errors in evaluating the estimated ratios are not
covered; verified bounds for them must be added. SVD-denoised estimates
are eligible only after their stated probability errors and normalized
mass vectors have been independently certified.

\begin{corollary}[One-sample conditional-frequency certificate]
\label{II:U:cor:rank-two-sampling}
Fix $0<\delta<1$, the two prefix words, and $E$ before drawing $n$ independent source
words, or use an independent validation split. Define
\[
 \nu_*=\min_{u\in\mathcal W}P_p(u\Sigma^*)>0.
\]
For each query use the empirical frequency of $E$ among suffixes of
sampled words beginning with that prefix. At the two anchors use the
empirical stop/next-letter histogram. If
$n\nu_*\ge8\log(3M/\delta)$, then with probability at least $1-\delta$
all prefix counts are at least $n\nu_*/2$ and the errors satisfy
\begin{align}
 \epsilon_E&\le
 \sqrt{\frac{\log(6M/\delta)}{n\nu_*}},\notag\\
 \epsilon_m&\le
 2\sqrt{\frac{(d+1)\log2+\log(6/\delta)}{n\nu_*}}.
 \label{II:U:eq:rank-two-sample-radii}
\end{align}
Whenever these radii meet the hypotheses of
Theorem~\ref{II:U:thm:rank-two-rounding}, its conclusion holds on this event.
\end{corollary}
\begin{proof}
Multiplicative Chernoff and a union bound make the probability of any
prefix count below half its expectation at most
$M\exp(-n\nu_*/8)\le\delta/3$.
Conditional on its count, the suffixes for a fixed prefix are iid from
that residual law. Hoeffding with a union bound over the $M$ event
queries gives the first radius with failure probability at most
$\delta/3$. For a histogram with $d+1$ categories, taking a union bound
over its $2^{d+1}$ sign vectors gives
$P(\|\widehat m-m\|_1>t)\le2^{d+1}\exp(-kt^2/2)$ at count $k$.
Apply this to the two anchor histograms and use $k\ge n\nu_*/2$ to
obtain the second radius with total failure probability at most
$\delta/3$. Dependence between different prefix queries causes no
problem for these union bounds.
\end{proof}

If a required prefix count is zero, the empirical construction returns
\texttt{UNRESOLVED}. A numerical certificate also requires verified
lower bounds for prefix masses and observable separation; unchecked
plug-in estimates of these lower bounds are not covered.

\paragraph{Depth selection and its scope.}
In Theorem~\ref{II:U:thm:rank-two-monotone}, compactness gives
\[
 w_*=\min_{\substack{a\in\Sigma\\t\in[r_-,r_+]}}w_a(t)>0,\qquad
 \mu=\min_{t\in[r_-,r_+]}\min(w_{a_-}(t),w_{a_+}(t))>0.
\]
All required queries, including the extra successor letter, satisfy
$\nu_*\ge w_*\mu^L$. Also
$\Delta_E^{(L)}\ge\Delta_E^{(1)}>0$.
Put $A=d+\log((d+1)/\delta)$,
$a=\log(1/\gamma_*)$, and $c=\log(1/\mu)$.
For sufficiently large $n/A$, the declared choice
\[
 L=\left\lfloor\frac{\log(n/A)}{2a+c}\right\rfloor\ge1
\]
and the preceding results give
\[
 \UTV(p,\widetilde q_{L,b})
 =O_{\mathfrak m}\left((A/n)^{a/(2a+c)}+d\,2^{-b}\right).
\]
The constants $\mathfrak m$ include the source contraction constants,
stop margin, event separation, and $w_*$. The count and small-error
conditions must still hold; they do for sufficiently large $n$ with
these fixed constants. This is a completed-word TV statement.
It assumes the extremal letters, event, and parameter bounds needed
for choosing $L$ are supplied or certified separately. Reusing the
same data for their adaptive choice needs a further uniform event.
It proves neither empirical singular-subspace stability nor a
depth-independent $n^{-1/2}$ boundary-estimation theorem, and it gives
no unqualified KL or uniform rare-prefix prediction guarantee.

\subsection{Restoring an anchored initial state}
\label{II:U:sec:rank-two-anchored}
\begin{proposition}[Three-state anchoring with correct occupation weights]
\label{II:U:prop:rank-two-anchored}
For any clipped two-prefix model above, adjoin a state representing
$p=f_{t_0}$, start there, keep its exact stop probability $s(t_0)$,
and send each letter to the clipped barycentric mixture in $[l,h]$.
Let $\beta$ be the total row of incoming probabilities into the two
old states and retain their aggregate matrix $Q$. The occupation
vector of the resulting proper three-state automaton is
\[
 x^{\rm anch}=(1,v),\qquad v=\beta(I-Q)^{-1}.
\]
If the invariant limiting interval is $[l_*,h_*]$, let
$\delta=\max(l-l_*,h_*-h)$. The exact defect is
\begin{align}
 J^{\rm anch}
 &=L_1\sum_a w_a(t_0)
       \operatorname{dist}(F_a(t_0),[l,h])\notag\\
 &\quad+L_1\sum_{t\in\{l,h\}}v_t
       \sum_a w_a(t)\operatorname{dist}(F_a(t),[l,h]),
 \label{II:U:eq:rank-two-anchored-defect}
\end{align}
and
$\UTV(p,q^{\rm anch})\le J^{\rm anch}/2
\le L_1(1-s_*)\delta/(2s_*)$.
\end{proposition}
\begin{proof}
The new state has no incoming transitions and is visited once.
The old states satisfy $v=\beta+vQ$. All three stops are at least
$s_*$, so their total expected visits are at most $1/s_*$.
The residual calculation is the same disjoint-first-letter calculation
as in \eqref{II:U:eq:rank-two-clip-defect}, including the new row.
Each displacement is at most $\delta$, giving the bound.
\end{proof}

In general $v\ne x_L$ and $v\ne x_L-\alpha_L$; neither identity may be
substituted without proof. On the golden-ratio source, for every
$L\ge1$ the initial successors $2/5,3/5$ are inside $[1-a,a]$,
so the new-row defect is zero. Symmetry gives
$\beta=(1/4,1/4)$ and $v=(1/2,1/2)$.
Consequently $J^{\rm anch}=J_L^{\rm clip}/2$ and
\[
 \UTV(p,q_L^{\rm anch})\le\frac{z_L}{20(1+z_L)^2}.
\]
The factor improvement is a certificate for this construction, not an
optimality claim. A pure-prefix two-state model uses a mixed initial
vector; imposing the empty-prefix initial state adds the third state.


\subsection{Observable coordinates and a finite positive certificate}
\label{II:U:sec:observable-positive}
The results below distinguish a supplied positive basis, selection from a
supplied candidate pool, and discovery of such a pool. Throughout this
continuation, matrix $\|\cdot\|_\infty$ is the maximum absolute row sum,
$\|\cdot\|_{\max}$ is the maximum absolute entry, and vector norms have
their usual meanings. All positive realizations are row-normalized:
$\alpha\mathbf1=1$ and $\sum_a A_a\mathbf1+\tau=\mathbf1$.

\begin{theorem}[Observable normalization and finite positivity certificate]
\label{II:U:thm:observable-certificate}
Let $p(w)=\alpha_*A_{*,w}\tau_*$ have a minimal $R$-state positive
realization with $\tau_*\ge s_*\mathbf1>0$. Fix distinct core prefixes
$u_i$ and suffixes $v_j$, $1\le i,j\le R$, for which
\[
 H_{ij}=p(u_iv_j),\quad (M_a)_{ij}=p(u_iav_j),\quad
 h_j=p(v_j),\quad t_i=p(u_i)
\]
give an invertible $H$. Put $B_a=M_aH^{-1}$, $\beta=hH^{-1}$,
$q_i=p(u_i\Sigma^*)$, and $D_q=\Udiag(q_i)$. Then $q_i>0$ and
\[
 q=(I-\textstyle\sum_aB_a)^{-1}t,\qquad
 \widetilde B_a=D_q^{-1}B_aD_q,\quad
 \widetilde\beta=\beta D_q,\quad \widetilde t=D_q^{-1}t
\]
satisfy $\widetilde\beta\mathbf1=1$ and
$\sum_a\widetilde B_a\mathbf1+\widetilde t=\mathbf1$.
For any invertible $V$ with $V\mathbf1=\mathbf1$, define
\[
 \alpha=\widetilde\beta V^{-1},\qquad
 A_a=V\widetilde B_aV^{-1},\qquad \tau=V\widetilde t.
\]
If $\alpha\ge0$, every $A_a\ge0$, and $\tau\ge s\mathbf1$ for some
$s>0$, these finitely many inequalities certify an exact proper
$R$-state realization of $p$.
\end{theorem}
\begin{proof}
Let $P_{i,:}=\alpha_*A_{*,u_i}$ and let column $j$ of $C$ be
$A_{*,v_j}\tau_*$. Since $H=PC$ is invertible, both factors are
invertible. Thus $B_a=PA_{*,a}P^{-1}$, $\beta=\alpha_*P^{-1}$,
and $t=P\tau_*$. Row normalization and the stopping margin give
$(I-\sum_aA_{*,a})^{-1}\tau_*=\mathbf1$, hence $q=P\mathbf1$.
Every row of $P$ is nonzero and nonnegative, proving $q_i>0$.
The normalized matrix $\widetilde P=D_q^{-1}P$ is row-stochastic;
the displayed identities follow by similarity. A second similarity
by $V$ preserves every coefficient. Moreover $V^{-1}\mathbf1=\mathbf1$,
so the resulting initial vector and transition/stop rows sum to one.
The nonnegative aggregate has row sums at most $1-s$; hence it is
transient, every state terminates, and the word law has mass one.
\end{proof}

The certificate eliminates an infinite family of word-positivity tests
\emph{after} $V$ is obtained. It does not find $V$. Observable coordinates
can be signed; they are not hidden-state posterior probabilities.

\subsection{Contracting cones and finite trailing-word dictionaries}
\label{II:U:sec:trailing-word}
\begin{theorem}[Conditional finite-prefix construction under projective contraction]
\label{II:U:thm:trailing-word}
Suppose a supplied finite-dimensional proper cone $\mathcal C$ carries
positive linear letter maps and a linear map $x\mapsto f_x$ to
nonnegative summable word functions. Let $\ell(x)=\sum_w f_x(w)$ be
strictly positive on $\mathcal C\setminus\{0\}$; normalized states have
$\ell(x)=1$. Assume letter shifts and normalized updates agree with this
representation, the initial law is $p=f_{x_0}$, and every normalized
state has stopping probability at least $s>0$. Each nonzero letter map
sends $\mathcal C\setminus\{0\}$ into its interior, has Hilbert diameter
at most $\Delta<\infty$, and has projective contraction coefficient at
most $\gamma\in(0,1)$. Then for every $L\ge1$ there is a proper automaton
$q_L$, using at most $N_d(L)=\sum_{j=0}^Ld^j$ actual-prefix residuals, with
\begin{equation}
 \UTV(p,q_L)\le\frac{1-s}{s}
       \bigl(e^{\Delta\gamma^{L-1}}-1\bigr).
 \label{II:U:eq:trailing-tv}
\end{equation}
The same bound holds after keeping one representative of each vertex of
the convex hull of the residuals of length at most $L$.
\end{theorem}
\begin{proof}
Use a state for every positive-mass prefix of length at most $L$, initialized at
$\varepsilon$. At state $u$, retain the source's stopping and letter
probabilities. For $|u|<L$, route letter $a$ to $ua$. For
$u=u_1\cdots u_L$, route it to $u_2\cdots u_La$. The latter source
state and the true successor are obtained by the \emph{same} sequence
of $L$ normalized letter maps, applied to $x_0$ and $x_{u_1}$.
The first map bounds their distance by $\Delta$, and the remaining
$L-1$ maps contract it to $D=\Delta\gamma^{L-1}$.

To pass to probability laws, take cone-order bounds
$m y\le_{\mathcal C}x\le_{\mathcal C}M y$. Positivity implies
$m f_y\le f_x\le M f_y$. Normalization gives $m\le1\le M$;
from $M/m\le e^D$ one gets $M\le e^D$ and hence
$\|f_x-f_y\|_1\le2(e^D-1)$, by summing the positive part.
Limits give the same bound when order extrema are not attained.
No coordinate $\ell^1$ bound on an arbitrary cone is used.

Each row defect is at most $2(1-s)(e^D-1)$. The constructed automaton
stops with probability at least $s$ at every visit, so its occupation
sum is at most $1/s$. The occupation-flow certificate gives
\eqref{II:U:eq:trailing-tv}. For pruning, first merge duplicate residuals,
retain the distinct vertices, and express the initial residual and each
routed successor as convex mixtures of retained vertices. The same row
defect and stopping estimates apply. Keeping the empty-prefix state
instead is always permitted within the original $N_d(L)$ budget.
\end{proof}

The Hilbert contraction estimate is the classical Birkhoff theorem
\citep{gaubert2015}. For fixed $s,\Delta,\gamma$ and $d\ge2$, the state bound is
$N(\epsilon)=O(\epsilon^{-\log d/\log(1/\gamma)})$:
use $e^z-1\le ze^\Delta$ for $0\le z\le\Delta$ and choose
$L=O(\log(1/\epsilon))$. This is exponential enumeration in depth,
not an efficient rank-$R$ discovery theorem. Exact transition queries,
cone data, and contraction bounds are inputs to this construction.

\subsection{Relative Doeblin anchors and a corrected simplex inclusion}
\label{II:U:sec:doeblin-sandwich}
Here $\Delta(Z)$ denotes the convex hull of the rows of a matrix $Z$.
\begin{lemma}[Inverse-entry bound for a nearly diagonal stochastic matrix]
\label{II:U:lem:inverse-entry}
If $\Pi\ge0$, $\Pi\mathbf1=\mathbf1$, and $\Pi_{ii}\ge1-\eta$
with $0\le\eta<1/2$, then $\Pi$ is invertible,
$\Pi^{-1}\mathbf1=\mathbf1$, and
\[
 (\Pi^{-1})_{ij}\ge-\frac{\eta}{1-2\eta}.
\]
\end{lemma}
\begin{proof}
For $\eta>0$, write $\Pi=(1-\eta)I+\eta K$ with $K$ row-stochastic.
Expand the inverse in powers of $-\eta K/(1-\eta)$. Each entry of
$K^n$ belongs to $[0,1]$. Dropping the even nonnegative terms bounds
every entry below by the sum of the odd negative coefficients,
$-\eta/(1-2\eta)$. The case $\eta=0$ is immediate.
\end{proof}


\begin{theorem}[Conditional Doeblin simplex inclusion and strict margins]
\label{II:U:thm:doeblin-sandwich}
Let a proper $R$-state law have $\tau_*\ge s_*\mathbf1>0$,
$\alpha_{*,j}\ge\mu$, and
\[
 A_{*,a,ij}\ge\mu w_{a,i},\qquad
 w_a=A_{*,a}\mathbf1,\qquad 0<\mu<1/R.
\]
Set $a_\mu=1-R\mu$, $J_\mu=a_\mu I+\mu\mathbf1\mathbf1^\top$,
and $V_*=\widetilde P^{-1}$ in Theorem~\ref{II:U:thm:observable-certificate}.
Suppose a supplied true anchor matrix has the form
$U=\Pi J_\mu V_*$, where $\Pi$ satisfies
Lemma~\ref{II:U:lem:inverse-entry}. Write
$c_{\max}=\max_j\sum_i\Pi_{ij}$ and choose a known upper bound
$C\ge c_{\max}$. Define
\[
 m_- =\frac{\eta}{1-2\eta},\qquad
 m_+=\frac{\mu}{a_\mu C}.
\]
If $m_-\le m_+$, every $m\in[m_-,m_+]$ gives an exact proper positive
realization through
\[
 D_m=(1+Rm)I-m\mathbf1\mathbf1^\top,\quad
 V=D_mU=WV_*,\quad W=D_m\Pi J_\mu.
\]
In particular,
$\Delta(J_\mu)\subseteq\Delta(W)\subseteq\Delta(I)$ and
$W\tau_*\ge s_*\mathbf1$. If $m_-<m<m_+$, put
$\delta_m=\min(m-m_-,m_+-m)$. If also
$w_{\min}=\min_{a,i}w_{a,i}>0$, all initial and transition entries are
at least
\begin{equation}
 \gamma_* =\frac{\delta_m}{1+Rm}\min(1,w_{\min})>0.
 \label{II:U:eq:sandwich-margin}
\end{equation}
One may use $C=1+(R-1)\eta$. A sufficient nonempty-interval condition is
\begin{equation}
 a_\mu(R-1)\eta^2+[1-(R-2)\mu]\eta-\mu\le0.
 \label{II:U:eq:sandwich-feasibility}
\end{equation}
\end{theorem}
\begin{proof}
Let $\widetilde W=D_m\Pi$. The inverse formula and row sums give
\[
 \widetilde W^{-1}=\frac{\Pi^{-1}+m\mathbf1\mathbf1^\top}{1+Rm}\ge0.
\]
It is row-stochastic. Since $J_\mu=\widetilde W^{-1}W$, this proves
the inner inclusion. For the outer inclusion one must retain column
sums $c_j$:
\[
 \widetilde W_{ij}=(1+Rm)\Pi_{ij}-mc_j,\qquad
 W_{ij}=a_\mu\widetilde W_{ij}+\mu
       \ge\mu-a_\mu m C\ge0.
\]
Also $W\mathbf1=\mathbf1$. For every row $z$ satisfying
$z_j\ge\mu(z\mathbf1)$,
\[
 (zJ_\mu^{-1})_j=\frac{z_j-\mu(z\mathbf1)}{a_\mu}\ge0.
\]
Consequently
\[
 \alpha=(\alpha_*J_\mu^{-1})\widetilde W^{-1},\qquad
 A_a=W(A_{*,a}J_\mu^{-1})\widetilde W^{-1},\qquad
 \tau=W\tau_*
\]
are nonnegative and normalized, with the asserted stopping margin.
If $m$ is strictly inside the interval, entries of $\widetilde W^{-1}$
are at least $\delta_m/(1+Rm)$ and $W_{ij}\ge a_\mu C\delta_m>0$.
The vector $\alpha_*J_\mu^{-1}$ has mass one. The row sums of
$W A_{*,a}J_\mu^{-1}$ equal $(Ww_a)_i\ge w_{\min}$, which proves
\eqref{II:U:eq:sandwich-margin}. Finally $c_j\le1+(R-1)\eta$; substituting
this bound into $m_-\le m_+$ gives \eqref{II:U:eq:sandwich-feasibility}.
\end{proof}

The margin in \eqref{II:U:eq:sandwich-margin} uses the stochasticity of $W$;
it does not require a cross-row letter-probability ratio. The hypotheses
concern relative anchors inside $\Delta(J_\mu)$, not nearly pure hidden
states in $\Delta(I)$. They are additional structural assumptions.

\subsection{Selection from a supplied relatively separable prefix pool}
\label{II:U:sec:relative-spa}
\begin{theorem}[Conditional relative SPA selection, with corrected noise budget]
\label{II:U:thm:relative-spa}
Let $G=J_\mu V_*$ be invertible, and let a supplied pool of $M$ true
prefix-coordinate rows form $F=HG$, where $H\ge0$ and
$H\mathbf1=\mathbf1$. Suppose there are distinct witness rows $i_r$
with $H_{i_r,r}\ge1-\eta_0$. Let
$\max_i\|\widehat F_{i,:}-F_{i,:}\|_2\le e_F$ and define
\begin{align*}
 \sigma&=\sigma_R(G),\qquad \kappa=\|G\|_2/\sigma,
 & C_{\rm SPA}&=1+80\kappa^2,\\
 e&=e_F+\sqrt2\eta_0\|G\|_2,
 & \theta_R&=\min\{(2\sqrt{R-1})^{-1},1/4\}.
\end{align*}
For $R\ge2$, if $e<\sigma\theta_R/C_{\rm SPA}$, ordinary SPA on
$\widehat F^\top$ returns $R$ indices in $O(MR^2)$ exact-real arithmetic
operations. After a permutation, their \emph{true} relative coordinates
satisfy
\begin{equation}
 H_{\widehat i_r,r}\ge1-\eta,\qquad
 \eta=\frac{e_F+C_{\rm SPA}e}{\sigma}.
 \label{II:U:eq:spa-relative-budget}
\end{equation}
If $\eta<1/2$ and \eqref{II:U:eq:sandwich-feasibility} is strict, these true
selected anchors satisfy Theorem~\ref{II:U:thm:doeblin-sandwich}.
\end{theorem}
\begin{proof}
Replace the witness rows of $F$ by the corresponding rows of $G$,
leaving other rows unchanged, to obtain an exactly separable matrix
$F^0$. The replacement changes each row by at most
$\sqrt2\eta_0\|G\|_2$, since a simplex vector $h$ with
$h_r\ge1-\eta_0$ obeys $\|h-e_r\|_2\le\sqrt2\eta_0$.
Apply the squared-Euclidean specialization of
\citet[Theorem 3 and Remark 2]{gillis2014} to $(F^0)^\top$ and
$\widehat F^\top$, using the larger factor $C_{\rm SPA}$.
Its hypotheses allow signed generator coordinates. The selected
\emph{observed} rows are within $C_{\rm SPA}e$ of distinct rows of $G$.
Their true rows are within $e_F+C_{\rm SPA}e$. Multiplication by
$G^{-1}$ bounds each relative-coordinate error by
$(e_F+C_{\rm SPA}e)/\sigma$, proving \eqref{II:U:eq:spa-relative-budget}.
\end{proof}

This is an application of a classical selection theorem, not a new SPA
algorithm. Its cost starts with the pool matrix already available.
The pool size, generation cost, prefix masses, and uniform estimation
error are separate inputs. The estimated rows $\widehat F_{\widehat i_r,:}$
are not the unknown true rows $F_{\widehat i_r,:}$; replacing the latter
in $V=D_mU$ requires an additional numerical error certificate.

\begin{proposition}[Uniform observable coordinates for a fixed prefix pool]
\label{II:U:prop:pool-estimation}
Use a fixed core with $|u_i|,|v_j|\le L_0$ and
$\sigma_R(H)\ge\sigma_0>0$, and prefixes in a fixed
pool $\mathcal P\subseteq\Sigma^{\le L_0}$ with masses at least $p_*>0$.
From $n$ independent words, observe the first $2L_0$ letters together
with whether and when the word terminates by that depth. Estimate
$H$, the core prefix masses $q$, the pool prefix masses $q_w$, and
$h_{w,j}=p(wv_j)$ by counts. Put
\[
 K=R^2+R+MR+M,\quad
 t=\sqrt{\log(2K/\delta)/(2n)},\quad h_0=\sqrt R/\sigma_0.
\]
If $Rh_0t\le1/2$ and $t\le p_*/2$, then with probability at least
$1-\delta$, all estimates
\[
 \widehat f_w=(\widehat h_w/\widehat q_w)\widehat H^{-1}D_{\widehat q}
\]
exist and have row $\ell^2$ error at most
\begin{equation}
 e_F\le\frac{8Rh_0t}{p_*}+(2Rh_0^2+h_0)t.
 \label{II:U:eq:pool-estimation}
\end{equation}
\end{proposition}
\begin{proof}
A union bound and Hoeffding give error at most $t$ for all $K$ entries.
The Neumann bound gives $\|\widehat H^{-1}\|_\infty\le2h_0$ and
$\|\widehat H^{-1}-H^{-1}\|_\infty\le2Rh_0^2t$.
Both conditional rows $h_w/q_w$ and $\widehat h_w/\widehat q_w$ are
nonnegative with mass at most one, because the $v_j$ are distinct.
Each coordinate differs by at most $4t/p_*$. Telescope
$(h_w/q_w)H^{-1}D_q$; the conditional-row error contributes at most
$8Rh_0t/p_*$ in $\ell^1$, and the remaining product error contributes
at most $(2Rh_0^2+h_0)t$. The $\ell^2$ norm is no larger.
\end{proof}

The uniform event permits choosing indices by SPA on the same pool.
It does not permit choosing the pool or core after observing these
data without an additional uniform argument or sample split.

\subsection{Pathwise stability and a complete fixed-basis sampling construction}
\label{II:U:sec:fixed-basis-sampling}
The following construction uses exact rational arithmetic for intermediate
linear algebra and a specified dyadic output. It is a finite-bit algorithm,
but it is not a theorem about unrestricted $b$-bit floating-point LU.

\begin{proposition}[Pathwise perturbation and normalized initial estimates]
\label{II:U:prop:pathwise-stability}
Let $p,\widehat q$ be proper row-normalized positive automata on the same
$R$ states, with $\tau\ge s\mathbf1>0$ for the first automaton. Then
\begin{equation}
 \UTV(p,\widehat q)\le\frac12\|\widehat\alpha-\alpha\|_1
 +\frac1{2s}\left(\|\widehat\tau-\tau\|_\infty+
                   \sum_a\|\widehat A_a-A_a\|_\infty\right).
 \label{II:U:eq:pathwise-general}
\end{equation}
If $\alpha_i,A_{a,ij}\ge\gamma>0$, $R\ge2$, and an unnormalized
estimate $a^0$ satisfies $\|a^0-\alpha\|_\infty\le\delta_\alpha
\le\gamma/[2(R+1)]$, its normalization $\bar\alpha=a^0/(a^0\mathbf1)$
is positive and
\begin{equation}
 \|\bar\alpha-\alpha\|_\infty
 \le\frac{(R-1)\delta_\alpha}{1-R\delta_\alpha}\le\gamma/2.
 \label{II:U:eq:initial-normalization}
\end{equation}
For transition estimates with per-letter maximum-entry errors $e_a$,
$\max_a e_a\le\gamma/2$ and $R\sum_a e_a\le s/2$ suffice for
$\bar A_{a,ij}\ge\gamma/2$ and
$\bar\tau=\mathbf1-\sum_a\bar A_a\mathbf1\ge(s/2)\mathbf1$.
\end{proposition}
\begin{proof}
For each $w=a_1\cdots a_m$, telescope
\begin{align*}
 \widehat\alpha\widehat A_w\widehat\tau-\alpha A_w\tau
 ={}&(\widehat\alpha-\alpha)\widehat A_w\widehat\tau
       +\alpha A_w(\widehat\tau-\tau)\\
 &+\sum_{k=1}^m\alpha A_{a_1\cdots a_{k-1}}
       (\widehat A_{a_k}-A_{a_k})
       \widehat A_{a_{k+1}\cdots a_m}\widehat\tau.
\end{align*}
The suffix sum $\widehat g=\sum_v\widehat A_v\widehat\tau$ is at most
$\mathbf1$ componentwise: its finite partial sums are
$\mathbf1-\widehat T^{n+1}\mathbf1$. Equality requires termination
from every state, and is not needed here. The original occupation row
$\nu=\alpha(I-T)^{-1}$ has $\nu\mathbf1\le1/s$.
Take absolute values, apply nonnegative matrix majorants, and sum over
all prefixes, the differing letter, and suffixes. Tonelli and these
two bounds give \eqref{II:U:eq:pathwise-general}. This also follows from
Lemma~\ref{II:U:lem:stopped-coupling}.

Write $a^0=\alpha+e$. In coordinate $i$ the normalized error is
$[e_i(1-\alpha_i)-\alpha_i\sum_{j\ne i}e_j]/(1+\sum_j e_j)$.
Its numerator is at most $(R-1)\delta_\alpha$ in absolute value;
$R\gamma\le1$ proves the last inequality in
\eqref{II:U:eq:initial-normalization}. Positivity of $a^0$ and of its sum
is immediate. Finally each row's total transition error is at most
$R\sum_a e_a$; recomputing STOP as the remainder proves the last claim.
\end{proof}


\begin{theorem}[Fixed rational basis: iid estimation, exact computation, dyadic output]
\label{II:U:thm:fixed-basis-sampling}
Fix the core in Theorem~\ref{II:U:thm:observable-certificate}, with
$|u_i|,|v_j|\le L_0$, $\sigma_R(H)\ge\sigma_0>0$, and
$q_i\ge q_*>0$. A \emph{supplied rational} invertible $V$ satisfies
$V\mathbf1=\mathbf1$ and certifies the true law with
$\alpha_i,A_{a,ij}\ge\gamma>0$, $\tau_i\ge s>0$, where the
supplied lower bounds $\gamma,s$ are rational.
The core and $V$ are fixed independently of the sample (or obtained
from an independent training sample with these true margins certified).
For $0<\delta<1$, observe $n$ iid words censored at $D=2L_0+1$, retaining whether and
when termination occurs by $D$, including at $D$ itself.

Estimate $H,M_a,h$ by exact completed-word counts and $q$ by prefix
counts, each divided by $n$. All subsequent intermediate operations
in this theorem are exact rational operations. Define
\begin{align*}
 t_n&=c_0/\sqrt n,&
 c_0&=\sqrt{\tfrac12\log(4R(R+2)/\delta)},\\
 b_n&=c_1/\sqrt n,&
 c_1&=\sqrt d+\sqrt{2\log(2R^2/\delta)},\\
 h_0&=\sqrt R/\sigma_0,&v_-&=\|V^{-1}\|_\infty,
 \qquad v_+=\|V\|_\infty,\\
 K_\alpha&=v_-[(2R+1)h_0+2Rh_0^2],\\
 K_0&=\frac{2Rh_0}{q_*^2}
          +\frac{4R^2h_0^2+4Rh_0}{q_*},\\
 E_\alpha&=K_\alpha c_0,&
 E_A&=v_+v_-\left(\frac{2h_0R^2}{q_*}c_1+K_0c_0\right).
\end{align*}
It suffices to take
\begin{equation}
 n\ge\max\left\{(2Rh_0c_0)^2,(2c_0/q_*)^2,
 (4E_\alpha/\gamma)^2,
 \left(2E_A/\min(\gamma,s)\right)^2\right\}.
 \label{II:U:eq:fixed-basis-n}
\end{equation}
Check that $\widehat H$ is invertible and all $\widehat q_i>0$;
otherwise return \texttt{UNRESOLVED}. Form
\begin{align*}
 a^0&=\widehat h\widehat H^{-1}D_{\widehat q}V^{-1},
 &\bar\alpha&=a^0/(a^0\mathbf1),\\
 \bar A_a&=V D_{\widehat q}^{-1}\widehat M_a
                  \widehat H^{-1}D_{\widehat q}V^{-1},
 &\bar\tau&=\mathbf1-\sum_a\bar A_a\mathbf1.
\end{align*}
Before dividing by $a^0\mathbf1$, check that this sum is positive.
Check normalization and the margins
$\bar\alpha_i,\bar A_{a,ij}\ge\gamma/2$ and
$\bar\tau_i\ge s/2$ exactly; on failure return
\texttt{UNRESOLVED}. The supplied checking thresholds can be taken
rational. Otherwise, with $h_b=2^{-b}\le\gamma/4$, round
every transition entry down to a multiple of $h_b$ and assign the
remaining row mass exactly to STOP. Round the first $R-1$ entries of
$\bar\alpha$ down and assign the remaining mass to its last entry.
The resulting dyadic law $q_b$ has exactly $R$ positive states and,
with probability at least $1-\delta$, is returned and satisfies
\begin{equation}
 \UTV(p,q_b)\le
 \frac{e_\alpha}{1-e_\alpha}+(R-1)h_b
          +\frac{e_A+dRh_b}{s},\qquad
 e_\alpha=E_\alpha/\sqrt n,\quad e_A=E_A/\sqrt n.
 \label{II:U:eq:fixed-basis-tv}
\end{equation}
In particular the rate is
$O_{R,\sigma_0,q_*,V,s,\gamma}
(\sqrt{(d+\log(R/\delta))/n}+d2^{-b})$.
The intermediate rational algorithm has polynomial bit complexity
in $n,D,d,R,b$ and the rational input bit length of $V$;
the output uses $O(dR^2b)$ bits for its probabilities.
\end{theorem}
\begin{proof}
For fixed $(i,j)$, the events $W=u_iav_j$ over $a$ are disjoint.
Cauchy--Schwarz bounds the expected sum of absolute frequency errors
by $\sqrt{d/n}$. Changing one observation changes that sum by at most
$2/n$. McDiarmid and a union bound over $R^2$ pairs, together with
Hoeffding for the $R^2+2R$ base entries, give a common event of
probability at least $1-\delta$ on which
\begin{equation}
 \|\widehat H-H\|_{\max},\ \|\widehat h-h\|_\infty,
 \|\widehat q-q\|_\infty\le t_n,
 \qquad
 \sum_a\|\widehat M_a-M_a\|_\infty\le R^2b_n.
 \label{II:U:eq:joint-count-event}
\end{equation}
The last inequality sums over both row and column indices. A
per-pair union bound cannot replace this step by exchanging a sum
over letters with a maximum over rows.

Since $\|H^{-1}\|_\infty\le h_0$, the first two sample thresholds give
\begin{align*}
 \|\widehat H^{-1}\|_\infty&\le2h_0,&
 \|\widehat H^{-1}-H^{-1}\|_\infty&\le2Rh_0^2t_n,\\
 \|D_{\widehat q}^{-1}\|_\infty&\le2/q_*,&
 \|D_{\widehat q}^{-1}-D_q^{-1}\|_\infty&\le2t_n/q_*^2.
\end{align*}
Distinct suffixes imply $\|h\|_1,\|\widehat h\|_1\le1$.
For each fixed core prefix, the words $u_iav_j$ over $(a,j)$ are
distinct, so
$\sum_a\|M_a\|_\infty,\sum_a\|\widehat M_a\|_\infty\le R$.
Also $\|D_q\|_\infty,\|D_{\widehat q}\|_\infty\le1$.
Telescoping the three factors in $hH^{-1}D_q$ gives
$\|a^0-\alpha\|_1\le K_\alpha t_n=e_\alpha$.
Telescoping $D_q^{-1}M_aH^{-1}D_q$ and summing over $a$ gives
\[
 \sum_a\|\bar A_a-A_a\|_\infty
 \le v_+v_-\left[
 \frac{2h_0}{q_*}\sum_a\|\widehat M_a-M_a\|_\infty+K_0t_n\right]
 \le e_A.
\]
The remaining thresholds imply $e_\alpha\le\gamma/4$ and
$e_A\le\min(\gamma,s)/2$. Thus $a^0$ is positive, and its sum is
within $e_\alpha$ of one. Its normalized entries are at least
$(\gamma-e_\alpha)/(1+e_\alpha)\ge\gamma/2$;
$\|\bar\alpha-\alpha\|_1\le2e_\alpha/(1-e_\alpha)$.
The transition entries are at least $\gamma/2$ and the recomputed
stopping probabilities at least $s/2$. All exact checks therefore pass.

Downward rounding changes at most $dR$ entries in each transition row,
with total change at most $dRh_b$; STOP absorbs this mass. Initial TV
changes by at most $(R-1)h_b$. The original source has occupation
sum at most $1/s$. Apply Lemma~\ref{II:U:lem:stopped-coupling}, or
\eqref{II:U:eq:pathwise-general} with the row-remainder identity, to obtain
\eqref{II:U:eq:fixed-basis-tv}. The rounding restriction leaves initial
and transition entries at least $\gamma/4$ and never reduces STOP.

For complexity, a trie of the $O(dR^2)$ queried words and prefixes
permits counting in $O(nD+dR^2D)$ symbol operations. Counts have
$O(\log n)$ bits. Fraction-free elimination computes the rational
inverses; determinants of the integer matrices have polynomial bit
length in $R,\log n$ and the input bit length of $V$. The remaining
$O(dR^3)$ dense rational arithmetic operations, normalization, and
integer floor divisions also have polynomial bit cost. This proves
the stated polynomial bound; it does not assert a small uniform
degree or a fixed-word-size implementation. Dyadic probabilities
are stored as their integer numerators, with remainders determined
exactly, giving the output bound.
\end{proof}


\paragraph{Scope of the numerical theorem.}
The source itself need not have rational coefficients: empirical
counts and the supplied basis are rational. Exact rational
intermediates avoid an unsupported universal backward-error constant
for partially pivoted LU. The result concerns the mathematical law
encoded by the dyadic probabilities, not bitwise equivalence to a
floating-point executable. A rational certified basis with strict
margins is a substantive input. Neither its existence for every
positive realization nor its recovery from data is asserted here.
Combining noisy SPA-selected rows with basis construction needs
control of the \emph{estimated basis}, not just selection of indices.
All these distinctions persist even though the fixed-basis sampling
and output-rounding result is now complete.

\subsection{Estimated bases and discovery in a bounded-depth promise class}
\label{II:U:sec:learned-basis}
The following perturbation argument closes the estimated-basis gap
for a stated finite-dimensional source class. All matrix norms in
this subsection are maximum absolute row sums.

\begin{lemma}[A perturbed normalized basis preserves a strict positive realization]
\label{II:U:lem:learned-basis}
Suppose $V\mathbf1=\mathbf1$ realizes a proper law with
$\alpha_j,A_{a,ij}\ge\gamma>0$ and $\tau_i\ge s>0$.
Let $\widehat V=V+E$, $E\mathbf1=0$, and
$\rho=\|E\|_\infty\|V^{-1}\|_\infty<1$.
In the new basis the exact coefficients of the \emph{same} law satisfy
\begin{align*}
 \|\alpha'-\alpha\|_1&\le\rho/(1-\rho),\\
 \sum_a\|A'_a-A_a\|_\infty&\le 2R\rho/(1-\rho),\\
 \|\tau'-\tau\|_\infty&\le\rho.
\end{align*}
Thus $\rho\le\min\{1/2,\gamma/(8R),s/2\}$ suffices for
$\alpha'_j,A'_{a,ij}\ge\gamma/2$ and $\tau'_i\ge s/2$.
Also $\|\widehat V^{-1}\|_\infty\le2\|V^{-1}\|_\infty$.
\end{lemma}
\begin{proof}
Set $K=EV^{-1}$, so $\widehat V=(I+K)V$, $K\mathbf1=0$.
Then $\alpha'=\alpha(I+K)^{-1}$,
$A'_a=(I+K)A_a(I+K)^{-1}$ and $\tau'=(I+K)\tau$.
The Neumann bound gives the first and third estimates, and
\[
 A'_a-A_a=(KA_a-A_aK)(I+K)^{-1}.
\]
Nonnegativity and normalization give
$\sum_a\|A_a\|_\infty\le R$, proving the second estimate.
The row identities remain exact; the margins imply positivity and
termination. The inverse bound is the same Neumann estimate.
\end{proof}


For an estimated observable row $z$ whose true row $f$ satisfies
$f\mathbf1=1$, use the affine normalization
\[
 \mathcal N(z)=z+(1-z\mathbf1)e_1^\top.
\]
Then $\mathcal N(z)\mathbf1=1$ exactly and
$\|\mathcal N(z)-f\|_1\le2\|z-f\|_1$. This works for signed
observable coordinates and avoids an unjustified probabilistic
interpretation of those coordinates.

\begin{theorem}[Frequent-prefix discovery, estimated basis, and dyadic law]
\label{II:U:thm:bounded-pool-learning}
Fix $R\ge2$, $d$, $L_0$, a supplied core of depth at most $L_0$,
and positive rational margins. Consider the following \emph{promise
class}: the source has a minimal proper $R$-state realization with
the Doeblin hypotheses of Theorem~\ref{II:U:thm:doeblin-sandwich},
stop margin $s_*$, $w_{\min}>0$, and $0<\mu<1/R$;
$\sigma_R(H)\ge\sigma_0$, core prefix masses are at least $q_*$;
for every relative vertex there is a witness prefix of depth at most
$L_0$, mass at least $p_*$, and relative purity at least $1-\eta_0$.
The realization and witness prefixes need not be supplied.

Choose $0\le\eta<1/2$ satisfying the strict inequality
\eqref{II:U:eq:sandwich-feasibility}, and a rational $m\in(m_-,m_+)$
with $C=1+(R-1)\eta$. Set
\begin{align*}
 a_\mu&=1-R\mu,& h_0&=\sqrt R/\sigma_0,& D_0&=1+2(R-1)m,\\
 s_G&=a_\mu/\sqrt R,& S_G&=\sqrt R h_0,
 & C_0&=1+80(S_G/s_G)^2,\\
 v_-&=\frac{1+(R-2)\mu}{a_\mu(1-2\eta)},&
 v_+&=D_0h_0,&
 \gamma_0&=\frac{\min(m-m_-,m_+-m)}{1+Rm}\min(1,w_{\min}).
\end{align*}
Observe $n$ iid words censored at $2L_0+1$, including termination
flags. Put $M_0=N_d(L_0)$,
\[
 K=M_0(R+1)+R^2+2R,\qquad
 t=\sqrt{\log(4K/\delta)/(2n)},\qquad
 b_n=\sqrt{d/n}+\sqrt{2\log(2R^2/\delta)/n}.
\]
Use the frequent observed-prefix pool
$\widehat{\mathcal P}=\{w:|w|\le L_0,\widehat q_w\ge3p_*/4\}$.
Compute its rows by Proposition~\ref{II:U:prop:pool-estimation}, then apply
$\mathcal N$. The following explicit sufficient tests are imposed:
\begin{align*}
 Rh_0t&\le1/2,& t&\le\min(p_*/4,q_*/2),\\
 e_0&=16Rh_0t/p_*+(2Rh_0^2+h_0)t,& e_F&=2e_0,\\
 e&=e_F+\sqrt2\eta_0S_G,&
 e&<s_G\theta_R/C_0,\qquad
 (e_F+C_0e)/s_G\le\eta,\\
 e_V&=D_0e_F,&
 e_Vv_-&\le\min\{1/2,\gamma_0/(8R),s_*/2\}.
\end{align*}
Here $\theta_R$ is as in Theorem~\ref{II:U:thm:relative-spa}.
Run SPA, let $\widehat U$ be its selected normalized rows, and set
$\widehat V=D_m\widehat U$. For the empirical realization define
\begin{align*}
 \bar v_-&=2v_-,&\bar v_+&=v_++e_V,\\
 e_\alpha&=\bar v_-[(2R+1)h_0+2Rh_0^2]t,\\
 K_0&=2Rh_0/q_*^2+(4R^2h_0^2+4Rh_0)/q_*,\\
 e_A&=\bar v_+\bar v_-
       [2h_0R^2b_n/q_*+K_0t].
\end{align*}
Require $e_\alpha\le\gamma_0/8$,
$e_A\le\min(\gamma_0,s_*)/4$, and $h_b=2^{-b}\le\gamma_0/8$.
Use exact rational arithmetic, normalize the initial vector, verify
initial/transition entries at least $\gamma_0/4$ and STOP at least
$s_*/4$, and round downward as in Theorem~\ref{II:U:thm:fixed-basis-sampling}.
On a failed check return \texttt{UNRESOLVED}.
With probability at least $1-\delta$ all checks pass and the returned
proper law on $R$ states satisfies
\begin{equation}
 \UTV(p,q_b)\le\frac{e_\alpha}{1-e_\alpha}+(R-1)h_b
                    +\frac{2(e_A+dRh_b)}{s_*}.
 \label{II:U:eq:learned-basis-tv}
\end{equation}
For fixed structural margins and sufficiently small $\eta_0$, the rate is
\[
 O\!\left(\sqrt{\frac{d+L_0\log d+\log(R/\delta)}n}+d2^{-b}\right).
\]
Counting uses a trie; the observed pool has at most $1+nL_0$ members.
With this supplied core the algorithm has polynomial bit cost in
$n,L_0,d,R,b$ and the rational parameter bit lengths.
\end{theorem}
\begin{proof}
Apply Hoeffding simultaneously to every prefix of depth at most
$L_0$ and its $R$ completed-word suffix queries, and to the core
queries. The union bound uses $K$ even though only observed frequent
prefixes are processed. Combine it with the joint-letter event in
\eqref{II:U:eq:joint-count-event}; the two failure budgets are $\delta/2$.
Every witness is retained, and every retained prefix has true mass
at least $p_*/2$. Proposition~\ref{II:U:prop:pool-estimation}, with this
lower mass, and affine normalization give row $\ell^1$ error $e_F$.

In native coordinates $\widetilde P=D_q^{-1}P$ is stochastic.
Thus $\sigma_R(J_\mu\widetilde P^{-1})\ge a_\mu/\sqrt R=s_G$.
The suffix matrix $C$ in $H=PC$ has nonnegative rows of mass at most
one, because suffixes are distinct. Hence
$\|V_*\|_\infty=\|CH^{-1}D_q\|_\infty\le h_0$ and
$\|J_\mu V_*\|_2\le S_G$.
The conservative SPA tests therefore imply that the selected
\emph{true} rows have relative purity $1-\eta$.

Their basis $V=D_mU$ has margins $\gamma_0,s_*$ by
Theorem~\ref{II:U:thm:doeblin-sandwich}. Every true observable residual
row has $\ell^1$ norm at most $h_0$, giving $\|V\|_\infty\le v_+$.
The identities
\[
 V^{-1}=\widetilde P J_\mu^{-1}\Pi^{-1}D_m^{-1},\quad
 \|\Pi^{-1}\|_\infty\le(1-2\eta)^{-1},\quad
 \|D_m^{-1}\|_\infty=1
\]
give the claimed $v_-$ bound. Moreover
$\|\widehat V-V\|_\infty\le e_V$ and both bases send
$\mathbf1$ to $\mathbf1$. Lemma~\ref{II:U:lem:learned-basis} implies
that the true law in the empirical basis has margins
$\gamma_0/2,s_*/2$ and the displayed inverse/norm bounds.

Now apply the deterministic product estimates from the proof of
Theorem~\ref{II:U:thm:fixed-basis-sampling} on the same uniform event.
They hold for every selected basis obeying these norm bounds; no
sample-independence claim for $\widehat V$ is needed. The raw initial
and transition errors are bounded by $e_\alpha,e_A$. The small-error
tests imply all stipulated margins after initial normalization.
Downward rounding preserves properness. Pathwise perturbation with
true STOP at least $s_*/2$ proves \eqref{II:U:eq:learned-basis-tv}.
A change of exact basis itself changes no word probability; it is
used to prove positivity, not counted as an additional distribution error.

All retained prefixes occur in the observed trie. Query lookup,
matrix products and SPA require polynomially many operations in its
size and the stated parameters. Squared-norm SPA comparisons can be
performed exactly on rationals; no square roots are needed to execute
the pivots. Fraction-free algebra and dyadic floor operations give
the polynomial bit bound as in Theorem~\ref{II:U:thm:fixed-basis-sampling}.
\end{proof}

This closes frequent-prefix discovery and the numerical basis bridge
\emph{within the stated promise class and with a supplied core}.
It does not certify the latent class assumptions from finite data,
discover the core in uniformly polynomial time, or use empirical SVD.
The finite exact positivity checks certify the output automaton;
the approximation guarantee additionally uses the source-class promise.

\begin{proposition}[Bounded-depth observations cannot certify an unrestricted source]
\label{II:U:prop:censored-obstruction}
For every censoring depth $D$, the laws $\delta_{a^{D+1}}$ and
$\delta_{a^{D+2}}$ induce identical censored observations, including
termination flags through $D$, but their TV distance is one.
For $\epsilon<1/2$ and $\delta<1/2$, no procedure using only these
observations can approximate both laws within $\epsilon$ with
probability at least $1-\delta$.
\end{proposition}
\begin{proof}
Every observation is the same nonterminated $D$-letter string.
The two TV balls of radius $\epsilon$ are disjoint, by the triangle
inequality. An output distribution that is identical under the two
sources cannot give both disjoint events probability greater than
$1/2$. Arbitrarily many such censored observations do not help.
\end{proof}


\begin{proposition}[Fixed-block empirical singular-space certificate]
\label{II:U:prop:fixed-block-svd}
Let $K$ be an $M\times N$ block of completed-word probabilities,
with its query words fixed before sampling, $\Urank K=r$, and
$\sigma_r(K)\ge s_0>0$. From $n$ iid observations let $\widehat K$
be the frequency matrix and put
\[
 e=\sqrt{MN}\sqrt{\log(2MN/\delta)/(2n)}.
\]
With probability at least $1-\delta$,
$\|\widehat K-K\|_2\le e$. If $e<s_0/2$, the leading $r$ left
and right singular projectors obey
\[
 \|\widehat P-P\|_2,\ \|\widehat Q-Q\|_2
          \le e/(s_0-e),\qquad
 \|\widehat K_r-K\|_2\le2e.
\]
If $3e<s_0$, thresholding empirical singular values at $2e$ selects
exactly $r$. The population rank and gap remain assumptions.
\end{proposition}
\begin{proof}
Hoeffding and a union bound give entry error at most
$e/\sqrt{MN}$ and thus the operator bound via the Frobenius norm.
The variational singular-value inequalities give
$\widehat\sigma_r\ge s_0-e$ and
$\widehat\sigma_{r+1}\le e$.
Writing $\widehat K\widehat V_r=\widehat U_r\widehat\Sigma_r$
and multiplying by $I-P$ gives
$(I-P)\widehat U_r=(I-P)(\widehat K-K)\widehat V_r
\widehat\Sigma_r^{-1}$. The norm is at most $e/(s_0-e)$.
For two projectors of equal rank this is their operator-norm distance;
the transposed argument gives the right projector. Finally
$\|\widehat K_r-K\|\le\|\widehat K_r-\widehat K\|+e\le2e$.
\end{proof}

This controls a fixed finite block. It neither estimates an infinite
Hankel tail nor supplies a positive cone or justifies adaptive block
selection outside a declared uniform family. Validated numerical
SVD errors must be added when the computed decomposition is inexact.

\subsection{A bounded-coordinate rank-two parametric construction}
\label{II:U:sec:parametric-rank-two}
\begin{theorem}[Estimated maps and endpoint clipping on a supplied valid interval]
\label{II:U:thm:parametric-rank-two}
Suppose $f_t$, $t\in I$, is a supplied affine family of probability
laws on a compact interval, with $p=f_{t_0}$,
$\|f_t-f_u\|_1\le K_f|t-u|$, true STOP at least $s>0$, and
letter masses $w_a(t)$ and projective updates $F_a(t)\in I$.
Let estimated affine-fractional maps $\widehat F_a:I\to I$, with
no pole on $I$, be contractions of modulus
at most $\bar\gamma<1$, and let nonnegative affine estimated masses
$\widehat w_a$ satisfy uniformly
\[
 \max_a|\widehat F_a(t)-F_a(t)|\le e_F,
 \qquad \sum_a|\widehat w_a(t)-w_a(t)|\le\zeta<s.
\]
Starting from $[t_0,t_0]$, perform $J$ all-letter endpoint-hull
iterations with $\widehat F_a$. At the two resulting endpoints,
use masses $\widehat w_a$, STOP as the remainder, and route to the
barycentric representation of $\widehat F_a(t)$ clipped to the
current interval. If the interval is degenerate, use one state.
Round probabilities down with step $h_b$ and initial normalization
by a remainder. The proper law on at most two states satisfies
\begin{equation}
 \UTV(p,q_{J,b})\le
 \frac{\zeta+\frac12K_f(e_F+\bar\gamma^J\operatorname{diam}I)}{s-\zeta}
       +h_b+\frac{2dh_b}{s-\zeta}.
 \label{II:U:eq:parametric-rank-two}
\end{equation}
Consequently fixed-depth coordinate estimates with
$e_F,\zeta=O(n^{-1/2})$ and $J=O(\log n)$ give a parametric
$O(n^{-1/2}+d2^{-b})$ TV rate. The interval and its validity are inputs.
\end{theorem}
\begin{proof}
The endpoint operator, taking minimum and maximum over $t_0$ and
all endpoint images, is $\bar\gamma$-Lipschitz in the maximum norm.
For an affine-fractional map without a pole on $I$, extrema on an
interval occur at its endpoints. Thus the interval expands monotonically and the
distance from any next image to the current interval is at most
$\bar\gamma^J\operatorname{diam}I$.
Initial mixing represents $f_{t_0}$ exactly. In a row, the STOP/letter
mass discrepancies contribute at most $2\zeta$ in $\ell^1$.
Since $\sum_a\widehat w_a\le1-s+\zeta\le1$, replacing true successors
by clipped estimated successors contributes at most
$K_f(e_F+\bar\gamma^J\operatorname{diam}I)$.
Every constructed STOP is at least $s-\zeta$. The occupation bound
and Theorem~\ref{II:U:thm:flows} give the first fraction. Initial dyadic TV
is at most $h_b$; per-row dyadic TV is at most $2dh_b$.
Apply the stopped-coupling bound for the remaining terms.
\end{proof}

For fixed-depth observable matrices, a certified nonzero denominator
on $I$ makes their fractional maps locally Lipschitz in the entries;
the inverse estimates already proved in Section~\ref{II:U:sec:fixed-basis-sampling}
then give the stated parametric errors. This does not assert that
the unknown valid interval, denominator margin or contraction bound
can always be inferred from bounded-depth data. Exact limiting rays
need not be evaluated: finite endpoint iteration and clipping suffice.

\section{Projective error, word laws and intermittent filtering}
\label{II:M:sec:projective-certificates}
Relative-error certificates can avoid an explicit
smallest-probability denominator, but their hypotheses must themselves
be established. They neither construct a positive cone from a Hankel
rank bound nor turn an average word-law guarantee into uniform factual
recall. Natural logarithms are used throughout this section.

\subsection{A sharp projective-to-KL bridge}
For positive vectors of the same dimension define
\[
 d_H(x,y)=\log\max_i\frac{x_i}{y_i}
          -\log\min_i\frac{x_i}{y_i}.
\]
This is a distance between rays: independent positive rescaling of
$x$ and $y$ does not change it. On distributions with common zeros,
the same definition may be used on their common positive support.

\begin{lemma}[Non-expansivity on the positive orthant]
\label{II:M:lem:projective-nonexpansive}
If $A$ is a nonnegative rectangular matrix with no zero row, then
$d_H(Ax,Ay)\le d_H(x,y)$ for all strictly positive $x,y$.
Normalizing either output by its sum leaves the conclusion unchanged.
\end{lemma}
\begin{proof}
Put $m=\min_i x_i/y_i$ and $M=\max_i x_i/y_i$. Multiplication by $A$
preserves $my\le x\le My$, giving $mAy\le Ax\le MAy$.
The outputs are strictly positive by the no-zero-row assumption.
Their coordinate ratios thus lie in $[m,M]$.
\end{proof}


\begin{lemma}[Sharp KL envelope at fixed projective diameter]
\label{II:M:lem:sharp-projective-kl}
Let $p,q$ be probability vectors positive on a common finite support.
If $d_H(p,q)\le e$, then
\begin{align}
 D_{\rm KL}(p\|q)&\le\Psi(e)\le e^2/8,\label{II:M:eq:sharp-projective-kl}\\
 \Psi(e)&=\frac{e}{\exp(e)-1}-1
          -\log\!\left(\frac{e}{\exp(e)-1}\right)\quad(e>0),
 \qquad \Psi(0)=0.\nonumber
\end{align}
The bound $\Psi(e)$ is attained by distributions on two points for each
$e>0$. Its small-$e$ expansion and large-$e$ behavior are
\[
 \Psi(e)=\frac{e^2}{8}-\frac{e^4}{576}
                 +\frac{e^6}{25920}+O(e^8),\qquad
 \Psi(e)=e-1-\log e+o(1)\quad(e\to\infty).
\]
\end{lemma}
\begin{proof}
Let $u_i=\log(p_i/q_i)$ and, with $I\sim p$, put
$W=\max_i u_i-u_I\in[0,e]$. Since $\mathbb E_p e^{-u_I}=1$,
\[
 D_{\rm KL}(p\|q)=\log\mathbb E e^W-\mathbb EW.
\]
For a fixed mean $\mu\in[0,e]$, convexity gives
$\mathbb E e^W\le1+(\mu/e)(\exp(e)-1)$, with equality for an
endpoint-supported law. Maximizing the resulting scalar expression
over $\mu$ gives $\mu_*=1-e/(\exp(e)-1)$ and the stated $\Psi(e)$.
To see attainability, give $W\in\{0,e\}$ mean $\mu_*$, call its
probabilities $p_i$, set $b=\log\mathbb E e^W$, and define
$q_i=p_i\exp(W_i-b)$. These $q_i$ sum to one, are positive, and
$\log(p_i/q_i)=b-W_i$ has oscillation $e$.

For the quadratic bound let $K(t)=\log\mathbb E e^{tW}$.
Its second derivative is the variance of $W$ under an exponential
tilt. Any variable in $[0,e]$ has variance at most $e^2/4$, since
$\operatorname{Var}(W)\le\mathbb E(W-e/2)^2\le e^2/4$.
Integrating $K''$ with weight $1-t$ on $[0,1]$ yields
$K(1)-K(0)-K'(0)\le e^2/8$.
Expanding the displayed elementary formula gives the asserted series;
using $e/(\exp(e)-1)\sim e\exp(-e)$ gives the large-$e$ expression.
\end{proof}


\subsection{What a relative matrix certificate actually asserts}
\begin{proposition}[Sufficient relative-error certificate]
\label{II:M:prop:relative-projective-certificate}
Let $A,B$ be nonnegative matrices with the same support and no zero
rows. If the oscillation of $\log(A_{ij}/B_{ij})$ on their support is
at most $\eta$, then $\sup_{x>0}d_H(Ax,Bx)\le\eta$.
In particular, if $0\le\delta<1$ and
$|B_{ij}-A_{ij}|\le\delta A_{ij}$ at every entry, then
\[
 \sup_{x>0}d_H(Ax,Bx)
 \le\log\frac{1+\delta}{1-\delta}
 =2\delta+O(\delta^3).
\]
The coefficient-oscillation condition is not necessary.
\end{proposition}
\begin{proof}
The first premise gives constants $m,M>0$ with $mB\le A\le MB$
and $\log(M/m)\le\eta$. Multiply by $x>0$ and compare output
ratios. In the relative-error case,
$(1-\delta)A\le B\le(1+\delta)A$, which proves the stated bound.
For failure of necessity take
\[
 A=\begin{pmatrix}1&1\\1&1\end{pmatrix},\qquad
 B=\begin{pmatrix}1&2\\1&2\end{pmatrix}.
\]
Both outputs are on the ray generated by $(1,1)^\top$, so their
projective distance is zero for every positive input. The coefficient
log-ratio oscillation is nevertheless $\log2$.
\end{proof}


With only the no-zero-row assumption, even equality of supports is
not necessary for a finite uniform projective defect: the matrices
\[
 \begin{pmatrix}1&0\\1&0\end{pmatrix},\qquad
 \begin{pmatrix}0&1\\0&1\end{pmatrix}
\]
again produce the same ray for every positive input. The nonzero-column
condition in Lemma~\ref{II:M:lem:projective-column-limit} below explains why
the normalized probabilistic setting is more restrictive.

\paragraph{Additive errors and zero entries.}
For known common support, minimum active entry $m>0$ and additive error
$\varepsilon<m$, the relative certificate gives
$\eta\le\log((1+\varepsilon/m)/(1-\varepsilon/m))$.
Dependence on $\varepsilon/m$ is unavoidable in worst cases:
$A=\operatorname{diag}(m,1)$ and
$B=\operatorname{diag}(m+\varepsilon,1)$ have defect
$\log(1+\varepsilon/m)$. This is not a universal equality in order:
$B=cA$ has zero projective defect despite a nonzero additive error.
Support-preserving quantization, multiplicative error bounds or other
separately proved restrictions are substantive premises. Smoothing is
one possible modification, not a free removal of these premises.

\begin{lemma}[Column limits of a uniform projective certificate]
\label{II:M:lem:projective-column-limit}
Suppose $A,B$ have no zero rows and every column of both has positive
sum. If $d_H(Ax,Bx)\le\eta$ for every $x>0$, then corresponding
normalized columns have identical supports and projective distance at
most $\eta$ on that support. Their KL divergence is at most $\Psi(\eta)$.
\end{lemma}
\begin{proof}
Apply the premise to $x=e_j+\epsilon\mathbf1$ and normalize both
outputs. The ratios of two normalized positive vectors with projective
distance at most $\eta$ lie in $[\exp(-\eta),\exp(\eta)]$.
Taking $\epsilon\downarrow0$ shows that each normalized column is
bounded above and below by these factors times the other. Thus their
supports coincide. Cross-ratio inequalities pass to the limit on that
support, so its projective distance is at most $\eta$.
Lemma~\ref{II:M:lem:sharp-projective-kl} proves the KL assertion.
\end{proof}


\subsection{Finite-horizon word laws without a mixing premise}
Use column-state conventions in this subsection. Two models have the
same finite latent state set $\{1,\ldots,N\}$, normalized positive
initial states $x_0,\widehat x_0$, and nonnegative matrices
$T_a,\widehat T_a$ with no zero rows. Define
\[
 C_{aj}=\mathbf1^\top T_a e_j>0,\qquad
 \widehat C_{aj}=\mathbf1^\top\widehat T_a e_j>0,
 \qquad \sum_a C_{aj}=\sum_a\widehat C_{aj}=1.
\]
From latent state $j$, the joint probability of emitting $a$ and moving
to $i$ is $(T_a)_{ij}$. Write $P_T,\widehat P_T$ for the induced
probability laws on $\Sigma^T$. These are fixed-horizon laws, not a
single normalized law on all finite words of all lengths.

\begin{theorem}[Word-law KL from aligned positive-model certificates]
\label{II:M:thm:projective-word-kl}
Assume for every positive $x$ and every letter $a$ that
\[
 d_H(T_ax,\widehat T_ax)\le\eta,
 \qquad d_H(Cx,\widehat Cx)\le\eta_{\rm obs}.
\]
Let $e_0=d_H(x_0,\widehat x_0)$. Then for every integer $T\ge1$,
\begin{align}
 D_{\rm KL}(P_T\|\widehat P_T)
 &\le\Psi(e_0)+T\bigl(\Psi(\eta_{\rm obs})+\Psi(\eta)\bigr),
 \label{II:M:eq:word-kl-relative}\\
 \frac1T D_{\rm KL}(P_T\|\widehat P_T)
 &\le\frac{e_0^2}{8T}+\frac{\eta_{\rm obs}^2+\eta^2}{8}.
 \label{II:M:eq:word-kl-quadratic}
\end{align}
Neither strict letter contraction nor intermittent mixing is required.
\end{theorem}
\begin{proof}
The emission distribution at latent state $j$ is $C_{\cdot j}$.
By Lemma~\ref{II:M:lem:projective-column-limit}, its KL divergence from
$\widehat C_{\cdot j}$ is at most $\Psi(\eta_{\rm obs})$.
Given $(j,a)$, the next-state distributions are
\[
 R_{a,j}=\frac{T_ae_j}{C_{aj}},\qquad
 \widehat R_{a,j}=\frac{\widehat T_ae_j}{\widehat C_{aj}}.
\]
The same lemma bounds their KL divergence by $\Psi(\eta)$ and
ensures matching support. The joint one-step kernels $K_j(a,i)$ and
$\widehat K_j(a,i)$ therefore satisfy, by the finite conditional
probability identity,
\[
 D_{\rm KL}(K_j\|\widehat K_j)
 =D_{\rm KL}(C_{\cdot j}\|\widehat C_{\cdot j})
  +\sum_a C_{aj}D_{\rm KL}(R_{a,j}\|\widehat R_{a,j})
 \le\Psi(\eta_{\rm obs})+\Psi(\eta).
\]
Expanding the log-likelihood ratio of the joint latent-state and letter
trajectory gives its initial KL plus the sum of the conditional kernel
KL divergences. The initial KL is at most $\Psi(e_0)$, so the displayed
bound follows for the joint trajectory. Marginalizing the latent states
cannot increase KL: grouping terms with the same observed word and
applying the log-sum inequality proves this directly. Finally apply
$\Psi(e)\le e^2/8$ to all three terms.
\end{proof}


\paragraph{Scope of the stronger word bound.}
The latent coordinates are aligned and both models have the same state
set. The uniform matrix certificates are stronger information than a
small empirical Hankel residual. The theorem does not construct such
coordinates, certify an arbitrary rank reduction, or infer the premises
from finitely many observed inputs. The total KL bound grows with $T$;
its rate bound does not imply a horizon-independent TV bound, uniform
error after conditioning on rare queries, or protected-fact persistence.
Termination and a law on $\Sigma^*$ require a separately specified
stopping model; they are not obtained by summing these normalized
fixed-length laws over $T$.

\subsection{Intermittent contraction controls the filter state}
For the models above set
$F_a(x)=T_ax/(\mathbf1^\top T_ax)$ and define $\widehat F_a$
similarly. Filters are evaluated on the same observed history drawn
under $P$. Non-expansivity holds even for identity or permutation
updates. It prevents amplification of an existing discrepancy but does
not, by itself, prevent accumulation of newly injected errors.

\begin{proposition}[Block filtering and the intermediate steps]
\label{II:M:prop:intermittent-projective-filter}
Assume the uniform update defect $\eta$ above. Suppose a set of
length-$L$ words has reference filter contraction coefficient at most
$\kappa_*<1$, and its conditional probability, under $P$ after every
supported history, is at least $p_*>0$. Put
\[
 q=\sqrt{1-p_*(1-\kappa_*^2)}<1,\qquad
 r_m=q^m e_0+L\eta\frac{1-q^m}{1-q}.
\]
With $e_t=d_H(x_t,\widehat x_t)$ one has, for $0\le j<L$,
\[
 \bigl(\mathbb E_P e_{mL}^2\bigr)^{1/2}\le r_m,
 \qquad
 \bigl(\mathbb E_P e_{mL+j}^2\bigr)^{1/2}\le r_m+j\eta.
\]
For $T=ML$, a valid filter-based word-law bound, including every
intermediate step, is
\[
 \frac1T D_{\rm KL}(P_T\|\widehat P_T)
 \le\frac1{8T}\sum_{m=0}^{M-1}\sum_{j=0}^{L-1}
                  (r_m+j\eta+\eta_{\rm obs})^2.
\]
The convention $q^0=1$ also covers $q=0$.
\end{proposition}
\begin{proof}
For a realized block $w$, telescoping its approximate updates against
reference updates and using non-expansivity after each substitution
gives $e_{(m+1)L}\le\kappa(w)e_{mL}+L\eta$.
Here $\kappa(w)\le\kappa_*$ on the mixing event and at most one
otherwise. Conditioning on the past yields
$\mathbb E[\kappa(w)^2\mid\mathcal F_{mL}]\le q^2$.
The $L^2(P)$ triangle inequality therefore gives
\[
 \|e_{(m+1)L}\|_{L^2(P)}
 \le q\|e_{mL}\|_{L^2(P)}+L\eta.
\]
Iteration proves the first bound; at an intermediate step use
$e_{mL+j}\le e_{mL}+j\eta$. The two next-letter distributions have
projective distance at most $e_t+\eta_{\rm obs}$. Apply
Lemma~\ref{II:M:lem:sharp-projective-kl}, the $L^2$ triangle inequality,
and the chain rule for the observed word law to obtain the full sum.
No independence of different blocks was used.
\end{proof}


Endpoint second-moment bounds do not on their own justify dropping the
$L-1$ intermediate steps from a word-law sum; the preceding proof includes
them. In the normalized,
aligned-model setting, Theorem~\ref{II:M:thm:projective-word-kl} gives the
simpler word-law estimate without any mixing premise. Existence and
efficient verification of $L,p_*,\kappa_*$ remain additional obligations
when a uniform filter-error bound is required.


\subsection{The columnwise hypothesis actually needed for word laws}
\begin{corollary}[Columnwise word-law certificate]\label{II:M:rev13:cor:column-word}
Consider two aligned $N$-state models with nonnegative matrices
$T_a,\widehat T_a$ and strictly positive normalized initial states
$x_0,\widehat x_0$. Let
\[
 C_{aj}=\mathbf1^\top T_ae_j>0,\quad
 \widehat C_{aj}=\mathbf1^\top\widehat T_ae_j>0,\quad
 \sum_a C_{aj}=\sum_a\widehat C_{aj}=1,
\]
and set $R_{a,j}=T_ae_j/C_{aj}$ and
$\widehat R_{a,j}=\widehat T_ae_j/\widehat C_{aj}$.
Assume corresponding $R_{a,j},\widehat R_{a,j}$ have common support and
\[
 d_H(R_{a,j},\widehat R_{a,j})\le\eta,\qquad
 d_H(Ce_j,\widehat Ce_j)\le\eta_{\rm obs}
 \quad\text{for every }a,j.
\]
With $e_0=d_H(x_0,\widehat x_0)$ and $T\ge1$,
\[
 D_{\rm KL}(P_T\Vert\widehat P_T)
 \le\Psi(e_0)+T\bigl(\Psi(\eta_{\rm obs})+\Psi(\eta)\bigr).
\]
These columnwise hypotheses suffice without a uniform projective
comparison of $T_ax$ and $\widehat T_ax$ for every positive $x$.
Zero rows of $T_a$ are allowed in this corollary. This weakening concerns
word-law KL; it does not replace the uniform update-error premise used
in the intermittent filter theorem.
\end{corollary}
\begin{proof}
For latent state $j$ the joint letter/next-state kernel is
$K_j(a,i)=C_{aj}R_{a,j}(i)$. Common support and the sharp envelope give
\[
 D_{\rm KL}(K_j\Vert\widehat K_j)
 =D_{\rm KL}(C_{\cdot j}\Vert\widehat C_{\cdot j})
   +\sum_a C_{aj}D_{\rm KL}(R_{a,j}\Vert\widehat R_{a,j})
 \le\Psi(\eta_{\rm obs})+\Psi(\eta).
\]
The initial KL is at most $\Psi(e_0)$. Apply the chain rule to the
length-$T$ latent trajectory and then marginalize the latent states.
\end{proof}


\section{Counterexamples for positive acquisition and conditional bounds}

\begin{proposition}[Uniform Doeblin positivity prevents arbitrary hidden-state purification]
\label{II:U:prop:doeblin-purity-obstruction}
If $A_{a,ij}\ge\mu(A_a\mathbf1)_i$ for every letter, each
nonempty-prefix posterior has all coordinates at least $\mu$.
The same is true of a normalized positive left Perron vector of a
nonempty word operator. Thus a collection of such anchors with
diagonal entries at least $1-\eta$ requires $\eta\ge(R-1)\mu$.
\end{proposition}
\begin{proof}
For a posterior $\pi$, $(\pi A_a)_j\ge\mu\pi A_a\mathbf1$.
Normalize and apply this observation to the last letter. A Perron
state returns to itself after the word and obeys the same inequality.
The other $R-1$ coordinates sum to at least $(R-1)\mu$.
\end{proof}

In particular $\eta<\mu s_*/(4R)$ is incompatible with these anchors
for $R\ge2$. An empty-prefix anchor supplies at most one exception.
The relative coordinates in Theorem~\ref{II:U:thm:doeblin-sandwich} avoid
this contradiction by using the inner simplex $J_\mu$.

\begin{proposition}[Noncommuting rank three with only one Perron ray]
\label{II:U:prop:common-perron}
Let $\alpha=(1,0,0)$ and
\[
 A_0=\frac1{30}\begin{pmatrix}4&1&1\\1&3&1\\1&2&4\end{pmatrix},
 \qquad
 A_1=\frac1{30}\begin{pmatrix}2&3&1\\3&1&2\\1&2&3\end{pmatrix}.
\]
With $\tau=(3/5,19/30,17/30)^\top$, this is a proper rank-three
law with noncommuting strictly positive transitions. Every nonempty
word operator has the same normalized left Perron vector
$(1/3,1/3,1/3)$.
\end{proposition}
\begin{proof}
The rows normalize and all stopping probabilities are positive.
For the core $\{\varepsilon,0,1\}$ on both sides, direct rational
calculation gives $\det H=-11/50625000\ne0$. Hence the rank is three.
The $(1,2)$ entry of $A_0A_1-A_1A_0$ is $1/450$.
Both matrices have column sums $1/5$, so their products have constant
column sum $5^{-|w|}$. Strict positivity makes the uniform left
Perron vector unique for each nonempty word.
\end{proof}


\paragraph{Dilation without a valid outer inclusion.}
Take $R=2$, $\eta=1/100$, $\mu=1/4$, $s_*=1/2$,
\[
 A_0=\begin{pmatrix}1/10000&1/10000\\1/8&1/8\end{pmatrix},\quad
 A_1=\tfrac14\mathbf1\mathbf1^\top-A_0,\quad
 \Pi=\begin{pmatrix}1&0\\1/100&99/100\end{pmatrix}.
\]
Although $\eta<\mu s_*/(4R)=1/64$, setting $m=1/98$ and
$W=D_m\Pi$ gives
\[
 W=\frac1{9800}\begin{pmatrix}9899&-99\\-1&9801\end{pmatrix},
 \qquad (WA_0W^{-1}\mathbf1)_1=-\frac{113851}{49000000}<0.
\]
This disproves the claimed algebraic positivity implication even
when a nearly diagonal stochastic anchor matrix is supplied.

\paragraph{The missing column-sum factor is substantive.}
With $R=3$, $\mu=1/5$, $\eta=1/10$, $m=49/100$, and
\[
 \Pi=\begin{pmatrix}1&0&0\\1/10&9/10&0\\0&0&1\end{pmatrix},
 \qquad
 W=D_m\Pi J_\mu
   =\frac1{2500}\begin{pmatrix}2431&59&10\\208&2282&10\\-39&59&2480\end{pmatrix},
\]
the old interval $1/8<m<1/2$ holds but $W_{31}<0$.
The corrected upper endpoint is $5/11$. This is realizable with
actual one-letter anchors, not only abstract matrices: put $Z=\Pi J_\mu$,
let the columns of
\[
 Q=\begin{pmatrix}1/4&1/1000&1/1000\\
                  1/10000&1/4&1/1000\\
                  1/10000&1/1000&1/4\end{pmatrix}
\]
be $w_a$, and define $A_a=w_aZ_{a,:}$, $a=0,1,2$, with uniform
initial distribution and STOP as the remainder. Every one-letter
posterior is the corresponding row of $Z$, all Doeblin inequalities
hold, and $\tau=(187/250,7489/10000,7489/10000)^\top$.
The one-letter prefix/suffix core has rank three. Nevertheless
$(WA_0W^{-1}\mathbf1)_3=-94961/25000000$.
Choosing $m=1/4$ in the corrected interval makes all transformed
coefficients strictly positive for this same source.

\subsection{Verified interior two-state example}
For the earlier source
$A_0=\frac1{20}\left(\begin{smallmatrix}3&2\\3&3\end{smallmatrix}\right)$,
$A_1=\frac1{20}\left(\begin{smallmatrix}2&2\\2&3\end{smallmatrix}\right)$,
$\alpha=(1/10,9/10)$, and $\tau=(11/20,9/20)^\top$, take
\[
 U=\begin{pmatrix}30/59&29/59\\20/49&29/49\end{pmatrix},\quad
 D(4,1)=\begin{pmatrix}2&-1\\-4&5\end{pmatrix},\quad
 V=\frac1{2891}\begin{pmatrix}1760&1131\\20&2871\end{pmatrix}.
\]
Then
\[
 \alpha_V=(897/5800,4903/5800),\qquad
 \tau_V=(29539/57820,26059/57820)^\top,
\]
\[
 A_{V,0}=\frac1{100606800}
       \begin{pmatrix}24761923&2357717\\24727123&5420117\end{pmatrix},
 \quad
 A_{V,1}=\frac1{100606800}
       \begin{pmatrix}16461862&5627438\\16427062&8689838\end{pmatrix}.
\]
These exact matrices certify $\gamma=1/50$ and $s=9/20$.
For the separate $R=2$ sandwich example with
$\mu=1/4$, $\Pi=\left(\begin{smallmatrix}1&0\\1/100&99/100\end{smallmatrix}\right)$,
and $m=1/98$, the correct values are
\[
 W=\frac1{19600}\begin{pmatrix}14799&4801\\4899&14701\end{pmatrix},
\]
\[
 WA_0W^{-1}=\frac1{970200000000}
       \begin{pmatrix}29484656149&30074228951\\
                       90086011249&91887363851\end{pmatrix}>0,
\]
where $A_0$ is the rare-letter matrix in the preceding obstruction.

\subsection{Why repeated whole-future KL cuts cannot be added}
\begin{proposition}[Counterexample to the proposed occupation-weighted KL bound]
\label{II:U:prop:repeated-kl-cuts}
The proposed lower bound $W_0\log(K/N)$ from $K>N$ disjoint future
signatures and cumulative prefix occupations is false without
nonoverlapping information increments.
\end{proposition}
\begin{proof}
On $\{0,a,b,c\}$ let
\[
 p=\tfrac13(\delta_{a0^{100}a}+\delta_{b0^{100}b}+\delta_{c0^{100}c}).
\]
For each $k\in\{a,b,c\}$ use prefixes $k0^j$, $0\le j\le100$,
and signature event $E_k=\{0^mk:m\ge0\}$. The 101 constant-length
cuts are prefix-free, the signatures are disjoint, and their
conditional probabilities equal one. Thus the proposed $N=2$, $K=3$
bound is $101\log(3/2)>40.95$. The one-state law with STOP $1/103$,
letter-$0$ probability $100/103$, and each other letter probability
$2/309$ is also representable by two identical states. Its divergence is
\[
 3\log103+100\log(103/100)+\log(3/4)<16.58.
\]
This contradicts the proposed bound. Each cut separately admits a
conditional KL lower bound, but its full future contains information
already present at the other cuts.
\end{proof}

If there are $M$ such cuts, summing their separate valid inequalities
and dividing by $M$ is legitimate. Adding them without this factor
requires signatures measurable in successive disjoint prediction
blocks and a corresponding chain rule. Neither version establishes
a general spectral KL discontinuity or directed conductance theorem.
